% !TeX root = ../main.tex
% Add the above to each chapter to make compiling the PDF easier in some editors.

\chapter{Approach}\label{chapter:approach}
\section{Design Goals}
\label{sec:design-goals}

The overarching goal of the proposed approach is to automatically augment an
existing seed corpus with inputs that enable a fuzzer to trigger regressions
introduced by a software commit as quickly as possible. This is motivated by the
integration of regression fuzzing into continuous integration and delivery
(CI/CD) pipelines, where fuzzing campaigns must provide useful feedback within a
limited time budget. Rather than requiring the fuzzer to discover the path to the
changed code through mutation alone, we aim to provide seeds that already
exercise code affected by the commit, giving the campaign a head start.

This objective translates into two Goals (G) for the resulting corpus:

\begin{description}
  \item[G1 --- Affected-code coverage.] The corpus should achieve high coverage
    of the code affected by the commit, as code that is never executed
    cannot expose a regression.
  \item[G2 --- Minimal corpus size.] The corpus should retain this coverage with
    as few seeds as possible, since processing a large corpus introduces
    additional I/O and substantial execution-time overhead during
    fuzzing~\cite{herreraCorpusDistillationEffective2020}.
\end{description}

\noindent
We therefore aim to construct a corpus that \emph{maximizes commit-affected code
coverage while minimizing corpus size}.

\paragraph{Augmentation instead of generation from scratch.}
We deliberately formulate this as a corpus \emph{augmentation} problem. In a
\ac{CI/CD} setting, projects are fuzzed repeatedly throughout their development,
allowing the seed corpus to accumulate inputs that exercise increasingly large
portions of the program. Discarding this accumulated corpus for every commit
would lose previously acquired knowledge and force the seed-generation process to
rediscover already established input structures. Instead, our approach reuses the
existing corpus and focuses the generation effort on increasing its coverage of
code affected by the current commit. This follows the established practice of
retaining and reusing fuzzing corpora across
campaigns~\cite{liImprovingSeedQuality2025}.

\paragraph{LLM Agent-based seed generation.}
For the augmentation itself, we employ a \ac{LLM}. Previous
work has demonstrated that \acp{LLM} can generate effective targeted fuzzing inputs by
reasoning about source code and input
structure~\cite{liuEvaluatingLLMBasedRegression, zhuSeedSmithLLMDrivenSeed2026,
xuDirectedGreyboxFuzzing2025, xuISC4DGFEnhancingDirected2024,
fengFuzzingRandomnessReasoning2025}. This capability complements conventional
mutation-based fuzzing, whose mutations are largely unaware of program semantics
and may therefore struggle to satisfy program-specific gating conditions such as
magic values, format constraints, or dependent validation checks. By analyzing
the source code alongside existing seeds, an \ac{LLM} agent can autonomously reason about these
conditions and construct inputs specifically designed to reach previously
uncovered program paths. In addition, by providing feedback signals from the affected code coverage report, these constructed inputs mainly cover regression-relevant code. The agent thus provides targeted, semantically informed
seeds, while the greybox fuzzer remains responsible for high-throughput mutation
and exploration around these starting points.

% ===========================================================================
%  pipeline.tex -- Corpus flow of the solution pipeline (overview)
%
%  Required in the preamble:
%      \usepackage{tikz}
%      \usetikzlibrary{arrows.meta,positioning,calc}
%
%  Pull it into a chapter with:   % ===========================================================================
%  pipeline.tex -- Corpus flow of the solution pipeline (overview)
%
%  Required in the preamble:
%      \usepackage{tikz}
%      \usetikzlibrary{arrows.meta,positioning,calc}
%
%  Pull it into a chapter with:   % ===========================================================================
%  pipeline.tex -- Corpus flow of the solution pipeline (overview)
%
%  Required in the preamble:
%      \usepackage{tikz}
%      \usetikzlibrary{arrows.meta,positioning,calc}
%
%  Pull it into a chapter with:   \input{figures/pipeline}
%  Natural size: about 14.4cm wide, fits \textwidth of the thesis (14.7cm).
% ===========================================================================
\begin{figure}
  \centering
  % ---------------------------------------------------------- icon library --
  \tikzset{
    icl/.style={line width=0.55pt, line cap=round, line join=round},
    pics/ic corpus/.style={code={
      \draw[icl] (-0.30,0.06) rectangle (0.10,0.34);
      \draw[icl,fill=white] (-0.20,-0.06) rectangle (0.20,0.22);
      \draw[icl,fill=white] (-0.10,-0.18) rectangle (0.30,0.10);
      \draw[icl] (-0.02,-0.02) -- (0.22,-0.02); \draw[icl] (-0.02,-0.10) -- (0.16,-0.10);
    }},
    pics/ic code/.style={code={
      \draw[icl] (-0.24,-0.32) -- (-0.24,0.32) -- (0.08,0.32) -- (0.24,0.16) -- (0.24,-0.32) -- cycle;
      \draw[icl] (0.08,0.32) -- (0.08,0.16) -- (0.24,0.16);
      \draw[icl] (-0.13,0.02) -- (-0.18,-0.06) -- (-0.13,-0.14);
      \draw[icl] (0.05,0.02) -- (0.10,-0.06) -- (0.05,-0.14);
    }},
    pics/ic diff/.style={code={
      \draw[icl] (-0.24,-0.32) -- (-0.24,0.32) -- (0.08,0.32) -- (0.24,0.16) -- (0.24,-0.32) -- cycle;
      \draw[icl] (0.08,0.32) -- (0.08,0.16) -- (0.24,0.16);
      \draw[icl] (-0.16,0.04) -- (0.02,0.04);  \draw[icl] (-0.07,-0.05) -- (-0.07,0.13);
      \draw[icl] (-0.16,-0.16) -- (0.02,-0.16);
    }},
    pics/ic gear/.style={code={
      \foreach \a in {0,45,...,315}{\draw[icl] (\a:0.18) -- (\a:0.32);}
      \draw[icl] (0,0) circle (0.19); \draw[icl] (0,0) circle (0.07);
    }},
    pics/ic cone/.style={code={
      \fill (-0.28,0) circle (0.06);
      \draw[icl] (-0.28,0) -- (0.28,0.26); \draw[icl] (-0.28,0) -- (0.28,-0.26);
      \draw[icl] (-0.28,0) -- (0.28,0.02);
      \draw[icl,dashed,dash pattern=on 1pt off 1pt] (0.28,0.26) -- (0.28,-0.26);
    }},
    pics/ic grid/.style={code={
      \foreach \x in {0,1,2}{\foreach \y in {0,1,2}{
        \draw[icl] (-0.30+\x*0.20,-0.30+\y*0.20) rectangle (-0.14+\x*0.20,-0.14+\y*0.20);}}
      \fill (-0.30,0.10) rectangle (-0.14,0.26);
      \fill (-0.10,-0.30) rectangle (0.06,-0.14);
      \fill (0.10,-0.10) rectangle (0.26,0.06);
    }},
    pics/ic robot/.style={code={
      \draw[icl] (0,0.34) -- (0,0.22); \fill (0,0.36) circle (0.04);
      \draw[icl,rounded corners=0.06cm] (-0.26,-0.20) rectangle (0.26,0.22);
      \fill (-0.12,0.04) circle (0.05); \fill (0.12,0.04) circle (0.05);
      \draw[icl] (-0.10,-0.10) -- (0.10,-0.10);
    }},
    pics/ic bug/.style={code={
      \draw[icl] (0,0.02) ellipse (0.17 and 0.24);
      \draw[icl] (0,0.30) circle (0.09);
      \draw[icl] (-0.05,0.37) -- (-0.12,0.46); \draw[icl] (0.05,0.37) -- (0.12,0.46);
      \foreach \y in {0.14,0.0,-0.14}{
        \draw[icl] (-0.16,\y) -- (-0.32,\y+0.06); \draw[icl] (0.16,\y) -- (0.32,\y+0.06);}
    }},
  }
  % \ic[colour]{icon}: icon on its own line
  \newcommand{\ic}[2][black!65]{\tikz[baseline=-0.30cm]{\pic[scale=0.82,color=#1]{#2};}}
  \begin{tikzpicture}[
      font   = \sffamily\footnotesize,
      >      = {Stealth[round]},
      box/.style   = {draw, align=center, inner sep=3pt, minimum height=17mm,
                      text width=17mm, rounded corners=2pt},
      corp/.style  = {box, fill=orange!14, draw=orange!70!black, line width=1pt},
      op/.style    = {box, fill=blue!7,    draw=blue!55},
      art/.style   = {box, fill=black!4,   draw=black!45},
      agent/.style = {box, fill=violet!14, draw=violet!60!black, line width=1pt},
      flow/.style     = {->, draw=black!60, line width=0.7pt},
      lane/.style     = {draw=black!60, line width=0.7pt},
      corpflow/.style = {->, draw=orange!70!black, line width=1.5pt},
      halo/.style     = {preaction={draw=white, line width=3pt, -}},
      el/.style    = {font=\sffamily\scriptsize, inner sep=1.5pt, fill=white,
                      text=black!70},
      stepno/.style  = {circle, fill=black!65, text=white, inner sep=0.6pt,
                      minimum size=4.2mm, font=\sffamily\tiny\bfseries},
    ]
    % ------------------------------------------------ corpus spine (y = 0) --
    \node[corp] (c0)   at ( 1.0,0) {\ic[orange!70!black]{ic corpus}\\[0.6mm]$C_0$\\reference};
    \node[op]   (dist) at ( 3.5,0) {\ic[blue!55!black]{ic grid}\\[0.6mm]corpus\\distillation};
    \node[corp] (c2)   at ( 6.0,0) {\ic[orange!70!black]{ic corpus}\\[0.6mm]$C_1$\\distilled};
    \node[corp] (ck)   at ( 8.5,0) {\ic[orange!70!black]{ic corpus}\\[0.6mm]$C_k$\\working};
    \node[corp] (c3)   at (11.0,0) {\ic[orange!70!black]{ic corpus}\\[0.6mm]$C^{*}$\\augmented};
    \node[op]   (fuzz) at (13.5,0) {\ic[blue!55!black]{ic bug}\\[0.6mm]fuzzing\\campaign};

    \draw[corpflow] (c0) -- (dist);
    \draw[corpflow] (dist) -- (c2);
    \draw[corpflow] (c2) -- (ck);
    \draw[corpflow] (ck) -- (c3);
    \draw[corpflow] (c3) -- (fuzz);

    % ------------------------------------------------ analysis (top row) ----
    \node[op] (comp) at (1.0,2.95) {\ic[blue!55!black]{ic gear}\\[0.6mm]compile\\target};
    \node[op] (coi)  at (6.0,2.95) {\ic[blue!55!black]{ic cone}\\[0.6mm]CoI\\computation};
    \node[op]  (llvm) at (3.5,2.95) {\ic[blue!55!black]{ic gear}\\[0.6mm]LLVM\\compile};
    \node[art] (diff) at (8.5,2.95) {\ic[black!55]{ic diff}\\[0.6mm]changed\\lines};
    \draw[flow] (llvm) -- node[el, above=1pt, fill=none] {IR} (coi);
    \draw[flow] (diff) -- (coi);

    % binary lane: compile -> distill, working corpus, fuzzing
    \draw[lane] (comp.south) -- (1.0,1.35) -- (13.5,1.35);
    \node[el] at (2.25,1.35) {binary};
    \draw[flow] (3.25,1.35) -- (3.25,0.85);
    \draw[flow] (8.75,1.35) -- (8.75,0.85);
    \draw[flow] (13.5,1.35) -- (fuzz.north);

    % affected-lines lane: CoI -> distill, working corpus
    \draw[lane] (coi.south) -- (6.0,1.72);
    \draw[lane] (3.75,1.72) -- (8.25,1.72);
    \node[el] at (4.88,1.72) {affected lines};
    \draw[flow,halo] (3.75,1.72) -- (3.75,0.85);
    \draw[flow,halo] (8.25,1.72) -- (8.25,0.85);

    % ------------------------------------------------ LLM agent loop --------
    \node[agent] (llm) at (8.5,-2.45) {\ic[violet!60!black]{ic robot}\\[0.6mm]LLM\\agent};
    \draw[flow] (8.25,-0.85) -- (8.25,-1.60);
    \draw[corpflow] (8.75,-1.60) -- (8.75,-0.85)
      node[el, midway, right=1pt, text=orange!65!black] {new seeds};

    % codebase the agent can inspect
    \node[art] (code) at (5.45,-2.45) {\ic[black!55]{ic code}\\[0.6mm]source\\code};
    \draw[flow] (code) -- node[el, below=1pt, fill=none] {inspect} (llm);

    % ------------------------------------------------ step numbers ----------
    \node[stepno] at (comp.north west) {1a};
    \node[stepno] at (llvm.north west) {1b};
    \node[stepno] at (coi.north west)  {1b};
    \node[stepno] at (dist.north west) {2};
    \node[stepno] at (llm.north west)  {3};
    \node[stepno] at (fuzz.north west) {4};
  \end{tikzpicture}
  \caption[Artifact flow of TACCO]{Artifact flow of TACCO.}
  \label{fig:solution-pipeline}
\end{figure}
%  Natural size: about 14.4cm wide, fits \textwidth of the thesis (14.7cm).
% ===========================================================================
\begin{figure}
  \centering
  % ---------------------------------------------------------- icon library --
  \tikzset{
    icl/.style={line width=0.55pt, line cap=round, line join=round},
    pics/ic corpus/.style={code={
      \draw[icl] (-0.30,0.06) rectangle (0.10,0.34);
      \draw[icl,fill=white] (-0.20,-0.06) rectangle (0.20,0.22);
      \draw[icl,fill=white] (-0.10,-0.18) rectangle (0.30,0.10);
      \draw[icl] (-0.02,-0.02) -- (0.22,-0.02); \draw[icl] (-0.02,-0.10) -- (0.16,-0.10);
    }},
    pics/ic code/.style={code={
      \draw[icl] (-0.24,-0.32) -- (-0.24,0.32) -- (0.08,0.32) -- (0.24,0.16) -- (0.24,-0.32) -- cycle;
      \draw[icl] (0.08,0.32) -- (0.08,0.16) -- (0.24,0.16);
      \draw[icl] (-0.13,0.02) -- (-0.18,-0.06) -- (-0.13,-0.14);
      \draw[icl] (0.05,0.02) -- (0.10,-0.06) -- (0.05,-0.14);
    }},
    pics/ic diff/.style={code={
      \draw[icl] (-0.24,-0.32) -- (-0.24,0.32) -- (0.08,0.32) -- (0.24,0.16) -- (0.24,-0.32) -- cycle;
      \draw[icl] (0.08,0.32) -- (0.08,0.16) -- (0.24,0.16);
      \draw[icl] (-0.16,0.04) -- (0.02,0.04);  \draw[icl] (-0.07,-0.05) -- (-0.07,0.13);
      \draw[icl] (-0.16,-0.16) -- (0.02,-0.16);
    }},
    pics/ic gear/.style={code={
      \foreach \a in {0,45,...,315}{\draw[icl] (\a:0.18) -- (\a:0.32);}
      \draw[icl] (0,0) circle (0.19); \draw[icl] (0,0) circle (0.07);
    }},
    pics/ic cone/.style={code={
      \fill (-0.28,0) circle (0.06);
      \draw[icl] (-0.28,0) -- (0.28,0.26); \draw[icl] (-0.28,0) -- (0.28,-0.26);
      \draw[icl] (-0.28,0) -- (0.28,0.02);
      \draw[icl,dashed,dash pattern=on 1pt off 1pt] (0.28,0.26) -- (0.28,-0.26);
    }},
    pics/ic grid/.style={code={
      \foreach \x in {0,1,2}{\foreach \y in {0,1,2}{
        \draw[icl] (-0.30+\x*0.20,-0.30+\y*0.20) rectangle (-0.14+\x*0.20,-0.14+\y*0.20);}}
      \fill (-0.30,0.10) rectangle (-0.14,0.26);
      \fill (-0.10,-0.30) rectangle (0.06,-0.14);
      \fill (0.10,-0.10) rectangle (0.26,0.06);
    }},
    pics/ic robot/.style={code={
      \draw[icl] (0,0.34) -- (0,0.22); \fill (0,0.36) circle (0.04);
      \draw[icl,rounded corners=0.06cm] (-0.26,-0.20) rectangle (0.26,0.22);
      \fill (-0.12,0.04) circle (0.05); \fill (0.12,0.04) circle (0.05);
      \draw[icl] (-0.10,-0.10) -- (0.10,-0.10);
    }},
    pics/ic bug/.style={code={
      \draw[icl] (0,0.02) ellipse (0.17 and 0.24);
      \draw[icl] (0,0.30) circle (0.09);
      \draw[icl] (-0.05,0.37) -- (-0.12,0.46); \draw[icl] (0.05,0.37) -- (0.12,0.46);
      \foreach \y in {0.14,0.0,-0.14}{
        \draw[icl] (-0.16,\y) -- (-0.32,\y+0.06); \draw[icl] (0.16,\y) -- (0.32,\y+0.06);}
    }},
  }
  % \ic[colour]{icon}: icon on its own line
  \newcommand{\ic}[2][black!65]{\tikz[baseline=-0.30cm]{\pic[scale=0.82,color=#1]{#2};}}
  \begin{tikzpicture}[
      font   = \sffamily\footnotesize,
      >      = {Stealth[round]},
      box/.style   = {draw, align=center, inner sep=3pt, minimum height=17mm,
                      text width=17mm, rounded corners=2pt},
      corp/.style  = {box, fill=orange!14, draw=orange!70!black, line width=1pt},
      op/.style    = {box, fill=blue!7,    draw=blue!55},
      art/.style   = {box, fill=black!4,   draw=black!45},
      agent/.style = {box, fill=violet!14, draw=violet!60!black, line width=1pt},
      flow/.style     = {->, draw=black!60, line width=0.7pt},
      lane/.style     = {draw=black!60, line width=0.7pt},
      corpflow/.style = {->, draw=orange!70!black, line width=1.5pt},
      halo/.style     = {preaction={draw=white, line width=3pt, -}},
      el/.style    = {font=\sffamily\scriptsize, inner sep=1.5pt, fill=white,
                      text=black!70},
      stepno/.style  = {circle, fill=black!65, text=white, inner sep=0.6pt,
                      minimum size=4.2mm, font=\sffamily\tiny\bfseries},
    ]
    % ------------------------------------------------ corpus spine (y = 0) --
    \node[corp] (c0)   at ( 1.0,0) {\ic[orange!70!black]{ic corpus}\\[0.6mm]$C_0$\\reference};
    \node[op]   (dist) at ( 3.5,0) {\ic[blue!55!black]{ic grid}\\[0.6mm]corpus\\distillation};
    \node[corp] (c2)   at ( 6.0,0) {\ic[orange!70!black]{ic corpus}\\[0.6mm]$C_1$\\distilled};
    \node[corp] (ck)   at ( 8.5,0) {\ic[orange!70!black]{ic corpus}\\[0.6mm]$C_k$\\working};
    \node[corp] (c3)   at (11.0,0) {\ic[orange!70!black]{ic corpus}\\[0.6mm]$C^{*}$\\augmented};
    \node[op]   (fuzz) at (13.5,0) {\ic[blue!55!black]{ic bug}\\[0.6mm]fuzzing\\campaign};

    \draw[corpflow] (c0) -- (dist);
    \draw[corpflow] (dist) -- (c2);
    \draw[corpflow] (c2) -- (ck);
    \draw[corpflow] (ck) -- (c3);
    \draw[corpflow] (c3) -- (fuzz);

    % ------------------------------------------------ analysis (top row) ----
    \node[op] (comp) at (1.0,2.95) {\ic[blue!55!black]{ic gear}\\[0.6mm]compile\\target};
    \node[op] (coi)  at (6.0,2.95) {\ic[blue!55!black]{ic cone}\\[0.6mm]CoI\\computation};
    \node[op]  (llvm) at (3.5,2.95) {\ic[blue!55!black]{ic gear}\\[0.6mm]LLVM\\compile};
    \node[art] (diff) at (8.5,2.95) {\ic[black!55]{ic diff}\\[0.6mm]changed\\lines};
    \draw[flow] (llvm) -- node[el, above=1pt, fill=none] {IR} (coi);
    \draw[flow] (diff) -- (coi);

    % binary lane: compile -> distill, working corpus, fuzzing
    \draw[lane] (comp.south) -- (1.0,1.35) -- (13.5,1.35);
    \node[el] at (2.25,1.35) {binary};
    \draw[flow] (3.25,1.35) -- (3.25,0.85);
    \draw[flow] (8.75,1.35) -- (8.75,0.85);
    \draw[flow] (13.5,1.35) -- (fuzz.north);

    % affected-lines lane: CoI -> distill, working corpus
    \draw[lane] (coi.south) -- (6.0,1.72);
    \draw[lane] (3.75,1.72) -- (8.25,1.72);
    \node[el] at (4.88,1.72) {affected lines};
    \draw[flow,halo] (3.75,1.72) -- (3.75,0.85);
    \draw[flow,halo] (8.25,1.72) -- (8.25,0.85);

    % ------------------------------------------------ LLM agent loop --------
    \node[agent] (llm) at (8.5,-2.45) {\ic[violet!60!black]{ic robot}\\[0.6mm]LLM\\agent};
    \draw[flow] (8.25,-0.85) -- (8.25,-1.60);
    \draw[corpflow] (8.75,-1.60) -- (8.75,-0.85)
      node[el, midway, right=1pt, text=orange!65!black] {new seeds};

    % codebase the agent can inspect
    \node[art] (code) at (5.45,-2.45) {\ic[black!55]{ic code}\\[0.6mm]source\\code};
    \draw[flow] (code) -- node[el, below=1pt, fill=none] {inspect} (llm);

    % ------------------------------------------------ step numbers ----------
    \node[stepno] at (comp.north west) {1a};
    \node[stepno] at (llvm.north west) {1b};
    \node[stepno] at (coi.north west)  {1b};
    \node[stepno] at (dist.north west) {2};
    \node[stepno] at (llm.north west)  {3};
    \node[stepno] at (fuzz.north west) {4};
  \end{tikzpicture}
  \caption[Artifact flow of TACCO]{Artifact flow of TACCO.}
  \label{fig:solution-pipeline}
\end{figure}
%  Natural size: about 14.4cm wide, fits \textwidth of the thesis (14.7cm).
% ===========================================================================
\begin{figure}
  \centering
  % ---------------------------------------------------------- icon library --
  \tikzset{
    icl/.style={line width=0.55pt, line cap=round, line join=round},
    pics/ic corpus/.style={code={
      \draw[icl] (-0.30,0.06) rectangle (0.10,0.34);
      \draw[icl,fill=white] (-0.20,-0.06) rectangle (0.20,0.22);
      \draw[icl,fill=white] (-0.10,-0.18) rectangle (0.30,0.10);
      \draw[icl] (-0.02,-0.02) -- (0.22,-0.02); \draw[icl] (-0.02,-0.10) -- (0.16,-0.10);
    }},
    pics/ic code/.style={code={
      \draw[icl] (-0.24,-0.32) -- (-0.24,0.32) -- (0.08,0.32) -- (0.24,0.16) -- (0.24,-0.32) -- cycle;
      \draw[icl] (0.08,0.32) -- (0.08,0.16) -- (0.24,0.16);
      \draw[icl] (-0.13,0.02) -- (-0.18,-0.06) -- (-0.13,-0.14);
      \draw[icl] (0.05,0.02) -- (0.10,-0.06) -- (0.05,-0.14);
    }},
    pics/ic diff/.style={code={
      \draw[icl] (-0.24,-0.32) -- (-0.24,0.32) -- (0.08,0.32) -- (0.24,0.16) -- (0.24,-0.32) -- cycle;
      \draw[icl] (0.08,0.32) -- (0.08,0.16) -- (0.24,0.16);
      \draw[icl] (-0.16,0.04) -- (0.02,0.04);  \draw[icl] (-0.07,-0.05) -- (-0.07,0.13);
      \draw[icl] (-0.16,-0.16) -- (0.02,-0.16);
    }},
    pics/ic gear/.style={code={
      \foreach \a in {0,45,...,315}{\draw[icl] (\a:0.18) -- (\a:0.32);}
      \draw[icl] (0,0) circle (0.19); \draw[icl] (0,0) circle (0.07);
    }},
    pics/ic cone/.style={code={
      \fill (-0.28,0) circle (0.06);
      \draw[icl] (-0.28,0) -- (0.28,0.26); \draw[icl] (-0.28,0) -- (0.28,-0.26);
      \draw[icl] (-0.28,0) -- (0.28,0.02);
      \draw[icl,dashed,dash pattern=on 1pt off 1pt] (0.28,0.26) -- (0.28,-0.26);
    }},
    pics/ic grid/.style={code={
      \foreach \x in {0,1,2}{\foreach \y in {0,1,2}{
        \draw[icl] (-0.30+\x*0.20,-0.30+\y*0.20) rectangle (-0.14+\x*0.20,-0.14+\y*0.20);}}
      \fill (-0.30,0.10) rectangle (-0.14,0.26);
      \fill (-0.10,-0.30) rectangle (0.06,-0.14);
      \fill (0.10,-0.10) rectangle (0.26,0.06);
    }},
    pics/ic robot/.style={code={
      \draw[icl] (0,0.34) -- (0,0.22); \fill (0,0.36) circle (0.04);
      \draw[icl,rounded corners=0.06cm] (-0.26,-0.20) rectangle (0.26,0.22);
      \fill (-0.12,0.04) circle (0.05); \fill (0.12,0.04) circle (0.05);
      \draw[icl] (-0.10,-0.10) -- (0.10,-0.10);
    }},
    pics/ic bug/.style={code={
      \draw[icl] (0,0.02) ellipse (0.17 and 0.24);
      \draw[icl] (0,0.30) circle (0.09);
      \draw[icl] (-0.05,0.37) -- (-0.12,0.46); \draw[icl] (0.05,0.37) -- (0.12,0.46);
      \foreach \y in {0.14,0.0,-0.14}{
        \draw[icl] (-0.16,\y) -- (-0.32,\y+0.06); \draw[icl] (0.16,\y) -- (0.32,\y+0.06);}
    }},
  }
  % \ic[colour]{icon}: icon on its own line
  \newcommand{\ic}[2][black!65]{\tikz[baseline=-0.30cm]{\pic[scale=0.82,color=#1]{#2};}}
  \begin{tikzpicture}[
      font   = \sffamily\footnotesize,
      >      = {Stealth[round]},
      box/.style   = {draw, align=center, inner sep=3pt, minimum height=17mm,
                      text width=17mm, rounded corners=2pt},
      corp/.style  = {box, fill=orange!14, draw=orange!70!black, line width=1pt},
      op/.style    = {box, fill=blue!7,    draw=blue!55},
      art/.style   = {box, fill=black!4,   draw=black!45},
      agent/.style = {box, fill=violet!14, draw=violet!60!black, line width=1pt},
      flow/.style     = {->, draw=black!60, line width=0.7pt},
      lane/.style     = {draw=black!60, line width=0.7pt},
      corpflow/.style = {->, draw=orange!70!black, line width=1.5pt},
      halo/.style     = {preaction={draw=white, line width=3pt, -}},
      el/.style    = {font=\sffamily\scriptsize, inner sep=1.5pt, fill=white,
                      text=black!70},
      stepno/.style  = {circle, fill=black!65, text=white, inner sep=0.6pt,
                      minimum size=4.2mm, font=\sffamily\tiny\bfseries},
    ]
    % ------------------------------------------------ corpus spine (y = 0) --
    \node[corp] (c0)   at ( 1.0,0) {\ic[orange!70!black]{ic corpus}\\[0.6mm]$C_0$\\reference};
    \node[op]   (dist) at ( 3.5,0) {\ic[blue!55!black]{ic grid}\\[0.6mm]corpus\\distillation};
    \node[corp] (c2)   at ( 6.0,0) {\ic[orange!70!black]{ic corpus}\\[0.6mm]$C_1$\\distilled};
    \node[corp] (ck)   at ( 8.5,0) {\ic[orange!70!black]{ic corpus}\\[0.6mm]$C_k$\\working};
    \node[corp] (c3)   at (11.0,0) {\ic[orange!70!black]{ic corpus}\\[0.6mm]$C^{*}$\\augmented};
    \node[op]   (fuzz) at (13.5,0) {\ic[blue!55!black]{ic bug}\\[0.6mm]fuzzing\\campaign};

    \draw[corpflow] (c0) -- (dist);
    \draw[corpflow] (dist) -- (c2);
    \draw[corpflow] (c2) -- (ck);
    \draw[corpflow] (ck) -- (c3);
    \draw[corpflow] (c3) -- (fuzz);

    % ------------------------------------------------ analysis (top row) ----
    \node[op] (comp) at (1.0,2.95) {\ic[blue!55!black]{ic gear}\\[0.6mm]compile\\target};
    \node[op] (coi)  at (6.0,2.95) {\ic[blue!55!black]{ic cone}\\[0.6mm]CoI\\computation};
    \node[op]  (llvm) at (3.5,2.95) {\ic[blue!55!black]{ic gear}\\[0.6mm]LLVM\\compile};
    \node[art] (diff) at (8.5,2.95) {\ic[black!55]{ic diff}\\[0.6mm]changed\\lines};
    \draw[flow] (llvm) -- node[el, above=1pt, fill=none] {IR} (coi);
    \draw[flow] (diff) -- (coi);

    % binary lane: compile -> distill, working corpus, fuzzing
    \draw[lane] (comp.south) -- (1.0,1.35) -- (13.5,1.35);
    \node[el] at (2.25,1.35) {binary};
    \draw[flow] (3.25,1.35) -- (3.25,0.85);
    \draw[flow] (8.75,1.35) -- (8.75,0.85);
    \draw[flow] (13.5,1.35) -- (fuzz.north);

    % affected-lines lane: CoI -> distill, working corpus
    \draw[lane] (coi.south) -- (6.0,1.72);
    \draw[lane] (3.75,1.72) -- (8.25,1.72);
    \node[el] at (4.88,1.72) {affected lines};
    \draw[flow,halo] (3.75,1.72) -- (3.75,0.85);
    \draw[flow,halo] (8.25,1.72) -- (8.25,0.85);

    % ------------------------------------------------ LLM agent loop --------
    \node[agent] (llm) at (8.5,-2.45) {\ic[violet!60!black]{ic robot}\\[0.6mm]LLM\\agent};
    \draw[flow] (8.25,-0.85) -- (8.25,-1.60);
    \draw[corpflow] (8.75,-1.60) -- (8.75,-0.85)
      node[el, midway, right=1pt, text=orange!65!black] {new seeds};

    % codebase the agent can inspect
    \node[art] (code) at (5.45,-2.45) {\ic[black!55]{ic code}\\[0.6mm]source\\code};
    \draw[flow] (code) -- node[el, below=1pt, fill=none] {inspect} (llm);

    % ------------------------------------------------ step numbers ----------
    \node[stepno] at (comp.north west) {1a};
    \node[stepno] at (llvm.north west) {1b};
    \node[stepno] at (coi.north west)  {1b};
    \node[stepno] at (dist.north west) {2};
    \node[stepno] at (llm.north west)  {3};
    \node[stepno] at (fuzz.north west) {4};
  \end{tikzpicture}
  \caption[Artifact flow of TACCO]{Artifact flow of TACCO.}
  \label{fig:solution-pipeline}
\end{figure}

\paragraph{Resulting pipeline.}
Based on these goals, we design the pipeline shown in
\cref{fig:solution-pipeline}, which proceeds in three stages:

\begin{enumerate}
  \item \textbf{Affected Code identification.} The code affected by the commit is
    identified through static code analysis. (G1+2)
  \item \textbf{Corpus distillation.} The existing corpus is reduced to a small
    set of seeds that efficiently covers this code (G2).
  \item \textbf{LLM-based augmentation.} An \ac{LLM} agent analyzes the program and
    the distilled corpus and generates additional seeds that increase
    affected-code coverage (G1).
\end{enumerate}

\noindent
The resulting corpus combines the accumulated knowledge of the existing corpus
with newly generated, commit-targeted inputs, and serves as the starting point
for the regression-fuzzing campaign.

\section{Computing the Commit Affected Code}
\label{sec:affected-code}

The first step of the pipeline is to identify the code affected by the commit, as illustrated by step~1b in \cref{fig:solution-pipeline}. The affected code comprises the parts of the program whose execution may be influenced by the changes introduced by the commit and can therefore potentially contain a regression resulting from these changes. Restricting the analysis to the lines directly modified by the commit would be insufficient, as a change can affect code execution beyond the site of change, and a resulting failure may manifest several lines away from the modified code. The broader motivation for considering affected code rather than only changed lines is discussed in \cref{chapter:background}.

We compute this affected code slice using CoIFuzz by Hundsdorfer~\cite{zotero-item-324}. As input, the pipeline first compiles the \ac{SUT} into LLVM bitcode. The resulting bitcode, together with the changed lines of the commit, is then provided to CoIFuzz. The tool then constructs a \ac{SVF} graph from the LLVM representation and traverses this graph to identify the code affected by the commit. The resulting slice is returned as a list of source files and line numbers.

The line-based representation of the affected slice also determines the coverage granularity used throughout the remainder of this thesis. Specifically, coverage is measured in terms of source lines, allowing the coverage achieved by individual seeds and corpora to be directly compared with the set of commit-affected lines identified by CoIFuzz. This \ac{CACC} is subsequently used as the optimization objective for corpus distillation and as the feedback signal during corpus augmentation.

\section{Distilling the Corpus}
\label{sec:distilling-corpus}

Before augmenting the corpus, we distill the reference corpus into a smaller set
of seeds that retains most of its \ac{CACC}. This step serves two purposes:

\begin{description}
  \item[Focused agent input.] It reduces the amount of material the \ac{LLM} agent
    must inspect, so that the agent sees fewer seeds with a higher information
    density per seed.
  \item[Lower fuzzing overhead.] It reduces the number of seeds the fuzzer must
    read, execute, and retain, and therefore the time spent on the initial
    corpus pass rather than on generating new inputs.
\end{description}

\noindent
Distillation is applied separately per project and fuzz target. It can also be
disabled entirely, which allows us to investigate the influence of different
starting corpora; this configuration is evaluated in
\cref{chapter:evaluation}.

\paragraph{Problem formulation.}
We formulate corpus distillation as a maximum coverage problem. Let
$S=\{s_1,\ldots,s_n\}$ denote the seeds remaining after the initial corpus
minimization described below, and let $A=\{a_1,\ldots,a_m\}$ denote the set of
affected source lines determined in \cref{sec:affected-code}.
For each seed $s_i$, let $C(s_i)\subseteq A$ denote the affected lines covered
by that seed. Given a corpus-size bound $k$, the objective is to select at most
$k$ seeds such that the number of distinct affected lines covered by the
selected seeds is maximized:

\begin{equation}
  \label{eq:max-coverage}
  \max_{S'\subseteq S}\ \left|\bigcup_{s\in S'} C(s)\right|
  \qquad\text{subject to}\qquad
  |S'|\leq k .
\end{equation}

\noindent
The constraint is an upper bound rather than an equality: $k$ is the
\emph{maximum permitted} corpus size, not the number of seeds to be selected,
so the optimization is free to return fewer seeds whenever this does not cost
coverage. Maximum coverage is NP-hard, and the greedy algorithm is optimal
among polynomial-time approximations under standard complexity
assumptions~\cite{feigeThresholdLnN1998}. The instances arising here are small
enough to be solved exactly and we solve each of them with \textit{PySAT}~\cite{PySATSATTechnology}.

\paragraph{Determining seed coverage.}
We first minimize the reference corpus with \texttt{afl-cmin} from
AFL++~\cite{fioraldiAFLCombiningIncremental2020}, which discards seeds that are
redundant with respect to edge coverage and thereby reduces the number of
candidates considered in the subsequent optimization. We then execute each
remaining seed individually on a coverage-instrumented build of the target and
intersect its source-line coverage with $A$. The resulting seed-to-line
coverage relation forms the input to the optimization problem. Because this
relation is measured rather than derived statically, seeds whose coverage is
input-independent but execution-dependent (for example due to timeouts) may be
recorded conservatively; we treat any line not observed as uncovered.

\paragraph{Constraint encoding.}
The seed-selection problem is encoded as a pseudo-Boolean optimization problem
over two families of binary variables. For each seed $s_i$ and affected line
$a_j$, let
\[
x_i =
\begin{cases}
1, & \text{if seed $s_i$ is selected},\\
0, & \text{otherwise},
\end{cases}
\qquad
y_j =
\begin{cases}
1, & \text{if affected line $a_j$ is covered},\\
0, & \text{otherwise}.
\end{cases}
\]

For every seed--line pair with $a_j \in C(s_i)$, selecting $s_i$ must imply that
$a_j$ is covered:
\begin{equation}
  \neg x_i \lor y_j .
  \label{eq:coverage-clause}
\end{equation}

Conversely, a line may be marked as covered only if at least one selected seed
covers it:
\begin{equation}
  \neg y_j
  \lor
  \bigvee_{i:\,a_j\in C(s_i)} x_i .
  \label{eq:coverage-definition}
\end{equation}

Together, \eqref{eq:coverage-clause} and \eqref{eq:coverage-definition}
establish that $y_j$ is true if and only if at least one selected seed covers
$a_j$. Only \eqref{eq:coverage-definition} is required for optimality, because
the objective rewards setting $y_j$ to true and therefore never leaves a
covered line unmarked; we include \eqref{eq:coverage-clause} so that the
variables carry their intended meaning in every model, not only in optimal
ones.

Finally, the corpus size is bounded by the cardinality constraint
\begin{equation}
  \sum_{i=1}^{n} x_i \leq k .
  \label{eq:size-constraint}
\end{equation}

\paragraph{Objective function.}
The primary objective is to maximize the number of distinct affected lines
covered by the selected corpus. As a secondary objective, solutions achieving
the same \ac{CACC} are preferred if they use fewer seeds. This lexicographic
objective is encoded as a single weighted objective:
\begin{equation}
  \max
  \left(
    (n+1)\sum_{j=1}^{m} y_j
    -
    \sum_{i=1}^{n} x_i
  \right).
  \label{eq:objective}
\end{equation}

\noindent
The weight $n+1$ strictly exceeds the largest possible value of
$\sum_i x_i$, so one additional covered line always outweighs any reduction in
corpus size.\footnote{Given \eqref{eq:size-constraint}, the weaker bound $k+1$
would suffice; we use $n+1$ so that the encoding remains valid independently of
$k$.} The objective therefore maximizes \ac{CACC} first and minimizes corpus
size only among solutions of equal coverage. The resulting corpus is
coverage-optimal under the given size bound and, among equally covering
solutions, minimal in the number of seeds. Ties beyond these two criteria are
broken by the solver; we fix the seed order and solver configuration so that
the selection is reproducible.

\paragraph{Choosing the corpus-size bound.}
Rather than committing to a single corpus-size bound, we solve
\eqref{eq:max-coverage} for a sequence of bounds, namely $k=1$ followed by
multiples of five, $k=5,10,15,20,\ldots$, up to the size of the minimized
corpus. Let $f(k)$ denote the optimal \ac{CACC} at bound $k$. Since the
feasible set grows with $k$, $f$ is non-decreasing, and the resulting coverage
curve exhibits diminishing returns. We select as the cutoff the smallest
evaluated bound $k_t$ whose relative gain over the following bound falls below
\SI{2}{\percent}, that is
\begin{equation}
  \frac{f(k_{t+1})-f(k_t)}{f(k_t)} < 0.02 .
  \label{eq:elbow}
\end{equation}
This threshold is a heuristic; it is fixed across all projects and was not
tuned per project.

The trade-off matters mainly for the subsequent \ac{LLM}-based augmentation. A
corpus of hundreds of seeds may provide marginally higher coverage, but it
exposes the agent to a large amount of input with little additional information
per seed. Applying \eqref{eq:elbow} reduces the reference corpus to between
\num{1} and \num{30} seeds while retaining between \SI{92.0}{\percent} and
\SI{100.0}{\percent} of the \ac{CACC} of the corresponding minimized corpus
across the evaluated projects (per-project means over the distillation runs;
see \cref{appendix:distill_elbow_curves}). Whether the sacrificed coverage can
subsequently be recovered, and whether the smaller corpus benefits the
augmentation process, is evaluated in \cref{chapter:evaluation}.

\paragraph{Role of the distilled corpus.}
The distilled corpus is not intended to be the final corpus. It provides a
compact, coverage-efficient starting point for the LLM agent, which
subsequently attempts to recover the uncovered affected code and to extend the
corpus with newly generated seeds.

\section{Augmenting the Corpus}
\label{sec:augmenting-corpus}

The final step of the pipeline augments the corpus with newly generated seeds that cover affected code not represented in the reference corpus. The goal of this step is not to generate arbitrary inputs, but to investigate the program and autonomously develop a strategy for constructing inputs that can reach the affected code.

This distinction motivates the use of an agent rather than a single-shot or fixed generation loop. Reaching a given affected line typically requires satisfying a chain of conditions that are themselves unrelated to the commit: magic bytes in a file header, length or offset fields that must remain consistent, checksums, or format-specific validation performed long before execution reaches the changed code. These gating conditions are not part of the affected slice and are therefore invisible in the coverage feedback. A generator presented only with the affected code and an aggregate coverage signal has no means of determining why its inputs are rejected: it observes that coverage does not increase, but not where execution is cut short. We initially implemented such an architecture, following the structure proposed by SeedMind~\cite{shiHarnessingLargeLanguage2024}, and observed exactly this failure mode. The model repeatedly produced inputs that were discarded during early parsing and received no information to correct them. Granting the model the ability to inspect the program itself, to locate the parser, read the validation logic, and examine existing seeds byte by byte, resolves this problem, because the information required to pass these gates is present in the source tree even though it is absent from the feedback signal. SeedSmith~\cite{zhuSeedSmithLLMDrivenSeed2026} likewise demonstrates that granting the model autonomy over its own investigation yields effective seeds, and that such autonomous exploration is a practical design for \ac{LLM}-driven seed generation.

To enable this behavior, the agent is provided with a workspace containing the artifacts it needs to reason about the program, a set of tools to explore them, and a single feedback tool to advance its corpus. \cref{fig:agent-loop} shows the resulting loop together with the artifacts and tools available at each stage. The remainder of this section first describes the environment in which the agent operates and then the procedure it follows within it.

% !TeX root = ../main.tex
% figures/agent_loop.tex  --  included via % !TeX root = ../main.tex
% figures/agent_loop.tex  --  included via % !TeX root = ../main.tex
% figures/agent_loop.tex  --  included via \input{figures/agent_loop.tex}
%
% Required in the preamble (already needed by figures/pipeline.tex):
%   \usepackage{tikz}
%   \usetikzlibrary{arrows.meta,positioning,calc}
%
% Drawing area: see \alk below (tuned for scrbook, fontsize=11pt, paper=a4).
% Do not scale this picture (scale=... shrinks the boxes but not the text).
% All text-bearing boxes size themselves to their content, so the figure
% survives a change of body/typewriter font.

\definecolor{alcorpus}{HTML}{D97706} % corpus spine (orange), match pipeline.tex
\definecolor{alagent}{HTML}{2563EB}  % agent / LLM and its tools
\definecolor{altool}{HTML}{475569}   % tooling and infrastructure
\definecolor{alprompt}{HTML}{7C3AED} % everything the LLM reads as a prompt
\definecolor{alcov}{HTML}{16A34A}    % [C] covered
\definecolor{alaff}{HTML}{DC2626}    % [A] affected
\colorlet{alold}{altool!55}          % older, superseded scripts

% ---- one knob: scales geometry and text together ---------------------------
% 1.125 -> main labels at 9pt = \footnotesize of scrbook (fontsize=11pt),
%          figure is about 14.3cm wide (\textwidth = 14.7cm), 11cm tall.
% 1.0   -> main labels at 8pt, about 12.7cm x 9.8cm.
\def\alk{1.125}
\pgfmathsetmacro\alsN{10*\alk}\pgfmathsetmacro\alsSm{9*\alk}
\pgfmathsetmacro\alsL{12*\alk}\pgfmathsetmacro\alsF{8*\alk}
\pgfmathsetmacro\alsS{7*\alk}\pgfmathsetmacro\alsT{5*\alk}
\def\alN {\fontsize{\alsN }{1.2\dimexpr\alsN  pt\relax}\selectfont}
\def\alSm{\fontsize{\alsSm}{1.2\dimexpr\alsSm pt\relax}\selectfont}
\def\alL {\fontsize{\alsL }{1.2\dimexpr\alsL  pt\relax}\selectfont}
\def\alF {\fontsize{\alsF }{1.2\dimexpr\alsF  pt\relax}\selectfont}
\def\alS {\fontsize{\alsS }{1.2\dimexpr\alsS  pt\relax}\selectfont}
\def\alT {\fontsize{\alsT }{1.2\dimexpr\alsT  pt\relax}\selectfont}

\pgfdeclarelayer{albg}               % card backgrounds drawn behind content
\ifdefined\altoolw\else\newlength{\altoolw}\fi
\settowidth{\altoolw}{\alS\ttfamily file\_type}

\begin{figure}[htbp]
  \centering
  \begin{tikzpicture}[
      scale=\alk,
      >={Latex[length=2.2mm,width=1.6mm]},
      font=\alF,
      %% ---------------- icons (TikZ native, no external fonts) --------- %%
      pics/folder/.style={code={
        \filldraw[draw=#1,fill=#1!14,line width=0.5pt,line join=round]
          (-0.21,-0.16) -- (-0.21,0.15) -- (-0.05,0.15) -- (0.00,0.09)
          -- (0.21,0.09) -- (0.21,-0.16) -- cycle;
        \draw[#1,line width=0.4pt] (-0.21,0.09) -- (0.00,0.09);}},
      pics/doc/.style={code={
        \filldraw[draw=#1,fill=#1!10,line width=0.5pt,line join=round]
          (-0.14,-0.19) -- (-0.14,0.19) -- (0.05,0.19) -- (0.14,0.10)
          -- (0.14,-0.19) -- cycle;
        \draw[#1,line width=0.4pt] (0.05,0.19) -- (0.05,0.10) -- (0.14,0.10);
        \draw[#1!75,line width=0.35pt]
          (-0.08,0.03) -- (0.08,0.03) (-0.08,-0.04) -- (0.08,-0.04)
          (-0.08,-0.11) -- (0.03,-0.11);}},
      pics/chip/.style={code={
        \foreach \y in {-0.08,0.00,0.08} {
          \draw[#1,line width=0.5pt] (-0.16,\y) -- (-0.23,\y);
          \draw[#1,line width=0.5pt] (0.16,\y) -- (0.23,\y);}
        \foreach \x in {-0.08,0.00,0.08} {
          \draw[#1,line width=0.5pt] (\x,0.16) -- (\x,0.23);
          \draw[#1,line width=0.5pt] (\x,-0.16) -- (\x,-0.23);}
        \filldraw[draw=#1,fill=#1!14,line width=0.6pt,rounded corners=1pt]
          (-0.16,-0.16) rectangle (0.16,0.16);
        \draw[#1,line width=0.5pt,rounded corners=0.5pt]
          (-0.07,-0.07) rectangle (0.07,0.07);}},
      pics/gear/.style={code={
        \foreach \a in {0,45,...,315} {
          \filldraw[draw=#1,fill=#1!20,line width=0.4pt]
            (\a-9:0.10) -- (\a-7:0.19) -- (\a+7:0.19) -- (\a+9:0.10) -- cycle;}
        \filldraw[draw=#1,fill=#1!14,line width=0.5pt] (0,0) circle (0.125);
        \filldraw[draw=#1,fill=white,line width=0.4pt] (0,0) circle (0.05);}},
      pics/todo/.style={code={
        \foreach \y/\done in {0.13/1,0.00/1,-0.13/0} {
          \draw[#1,line width=0.4pt,fill=white] (-0.20,\y-0.045) rectangle ++(0.09,0.09);
          \ifnum\done=1
            \draw[#1,line width=0.6pt] (-0.19,\y) -- (-0.16,\y-0.03) -- (-0.10,\y+0.05);
          \fi
          \draw[#1!60,line width=0.9pt,line cap=round] (-0.04,\y) -- (0.18,\y);}}},
      pics/check/.style={code={
        \draw[#1,line width=1.0pt,line cap=round,line join=round]
          (-0.08,0.00) -- (-0.02,-0.06) -- (0.09,0.07);}},
      %% ---------------- boxes ------------------------------------------ %%
      container/.style={
        draw=altool!50, dashed, rounded corners=3pt, line width=0.5pt},
      ctitle/.style={font=\alF\bfseries, text=altool},
      artifact/.style={
        draw=altool!60, fill=altool!6, rounded corners=2pt, line width=0.5pt,
        minimum width=\alk*4.0cm, inner sep=0pt},
      corpusartifact/.style={
        artifact, draw=alcorpus, fill=alcorpus!10, line width=0.7pt},
      alabel/.style={anchor=west, align=left, inner sep=0pt},
      agentbox/.style={
        draw=alagent, fill=alagent!8, rounded corners=3pt, line width=0.8pt,
        minimum width=\alk*2.8cm, minimum height=\alk*1.50cm, inner sep=0pt},
      toolbadge/.style={
        draw=alagent!70, fill=white, rounded corners=1.5pt, line width=0.45pt,
        font=\alS\ttfamily, text=alagent!80!black,
        minimum width=\altoolw+3pt, minimum height=\alk*0.29cm, inner sep=1pt},
      promptcard/.style={
        draw=alprompt, fill=alprompt!6, rounded corners=2pt, line width=0.8pt},
      prompttag/.style={
        fill=alprompt, text=white, font=\alT\bfseries\sffamily,
        rounded corners=1pt, inner sep=1.6pt, anchor=south west},
      finalbox/.style={
        draw=alcorpus, fill=alcorpus!22, rounded corners=2pt, line width=0.8pt,
        minimum width=\alk*4.0cm, minimum height=\alk*0.80cm, inner sep=0pt},
      tooltag/.style={
        fill=alagent, text=white, font=\alT\bfseries\sffamily,
        rounded corners=1pt, inner sep=1.6pt, anchor=south west},
      iotag/.style={
        fill=alcorpus, text=white, font=\alT\bfseries\sffamily,
        rounded corners=1pt, inner sep=1.6pt, anchor=south west},
      steprow/.style={
        draw=altool!45, fill=white, rounded corners=2pt, line width=0.5pt,
        minimum width=\alk*5.25cm, minimum height=\alk*0.80cm, inner sep=0pt},
      stepnum/.style={
        circle, draw=altool, fill=altool, text=white, inner sep=0pt,
        minimum size=\alk*0.46cm, font=\alT\bfseries\sffamily},
      checkbadge/.style={
        draw=alcov!70, fill=alcov!8, rounded corners=1.5pt, line width=0.4pt,
        font=\alT, text=alcov!70!black, minimum width=\alk*0.78cm,
        text height=\alk*0.38cm, text depth=\alk*0.06cm, inner xsep=2pt, inner ysep=1.5pt},
      icap/.style={font=\alT, text=altool, inner sep=0pt},
      gmark/.style={font=\alS\ttfamily\bfseries, inner sep=0pt},
      codeline/.style={line width=2.2pt, line cap=round, draw=altool!30},
      flow/.style={->, draw=altool, line width=0.7pt},
      mini/.style={->, draw=altool!80, line width=0.5pt,
        >={Latex[length=1.3mm,width=1.0mm]}},
      agentflow/.style={->, draw=alagent, line width=0.8pt},
      promptflow/.style={->, draw=alprompt, line width=1.0pt},
      corpusflow/.style={->, draw=alcorpus, line width=1.0pt},
      halo/.style={preaction={draw=white, line width=3.2pt, -}},
      lbl/.style={font=\alS, text=altool, inner sep=1.5pt},
      corpuslbl/.style={lbl, text=alcorpus},
    ]

    \pgfsetlayers{albg,main}

    %% ================================================================ %%
    %% Workspace  (container 0 .. 4.4  /  -3.35 .. 3.30)
    %% ================================================================ %%
    \draw[container] (0,-3.35) rectangle (4.4,3.22);
    \node[ctitle] at (2.2,3.03) {Workspace};

    %% --- generator scripts (one per cycle) --------------------------- %%
    \node[artifact,minimum height=\alk*1.00cm] (script) at (2.2,2.40) {};
    \pic[scale=0.85] at (0.50,2.64) {doc=alold};
    \node[alabel,text=altool!70] at (0.80,2.64) {\ttfamily seed\_gen\_cycle\_1.py};
    \pic[scale=0.85] at (0.50,2.18) {doc=altool};
    \node[alabel] at (0.80,2.18) {\ttfamily seed\_gen\_cycle\_N.py};

    %% --- annotated source (mini listing with gutter markers) -------- %%
    \node[artifact,minimum height=\alk*2.90cm] (src) at (2.2,0.25) {};
    \pic at (0.50,1.40) {folder=altool};
    \node[alabel] at (0.80,1.40)
      {\ttfamily src/};
    \filldraw[draw=altool!30,fill=white,rounded corners=1.5pt,line width=0.4pt]
      (0.35,1.12) rectangle (4.05,-0.42);
    \foreach \y/\m/\c/\x/\len in {
        0.95/[C]/alcov/1.05/2.3,
        0.65/[A]/alaff/1.35/2.3,
        0.35/[P]/alcov!55/1.35/1.5,
        0.05/[A]/alaff/1.65/1.9,
       -0.25/[.]/altool!45/1.05/1.1} {
      \node[gmark,text=\c] at (0.65,\y) {\m};
      \draw[codeline] (\x,\y) -- ++(\len,0);
    }
    \foreach \x/\y/\c/\m/\t in {
        0.40/-0.72/alaff/[A]/affected,  2.25/-0.72/alcov/[C]/covered,
        0.40/-0.98/alcov!55/[P]/partial, 2.25/-0.98/altool!45/[.]/uncovered}
      \node[icap,font=\alS,anchor=west] at (\x,\y)
        {\textcolor{\c}{\ttfamily\bfseries\m}\ \t};

    %% --- corpus: reference + generated seeds ------------------------ %%
    \node[corpusartifact,minimum height=\alk*1.40cm] (corpus) at (2.2,-2.25) {};
    \pic at (0.50,-1.85) {folder=alcorpus};
    \node[alabel] at (0.80,-1.85) {\ttfamily corpus/};
    \foreach \x in {0.65,0.95,1.25} \pic[scale=0.75] at (\x,-2.35) {doc=altool};
    \foreach \x in {2.35,2.65,2.95,3.25} \pic[scale=0.75] at (\x,-2.35) {doc=alcorpus};
    \node[icap] at (0.95,-2.72) {reference};
    \node[icap,text=alcorpus] at (2.80,-2.72) {generated};

    %% --- output ------------------------------------------------------ %%
    \node[finalbox] (final) at (2.2,-4.45) {};
    \node[iotag] at (final.north west) {OUTPUT};
    \pic[scale=1.3] at (0.60,-4.45) {folder=alcorpus};
    \node[alabel,font=\alF\bfseries,text=alcorpus!70!black]
      at (1.05,-4.45) {augmented corpus};
    \draw[corpusflow] (corpus.south) -- (final.north);
    \node[lbl,anchor=west] at (2.35,-3.65) {budget exhausted};

    %% --- input ------------------------------------------------------- %%
    \node[finalbox] (distilled) at (2.2,3.72) {};
    \node[iotag] at (distilled.north west) {INPUT};
    \pic[scale=1.3] at (0.60,3.72) {folder=alcorpus};
    \node[alabel,font=\alF\bfseries,text=alcorpus!70!black]
      at (1.05,3.72) {distilled corpus};
    \draw[corpusflow,line width=1.4pt] (distilled.west) -| (-0.30,-2.25)
      -- (corpus.west);

    %% ================================================================ %%
    %% Agent  (7.75 .. 10.75 / 0.70 .. 2.40) and its todo_list tool
    %% ================================================================ %%
    \node[agentbox] (agent) at (9.35,1.55) {};
    \pic[scale=1.4] at (8.40,1.55) {chip=alagent};
    \node[align=center,font=\alN\bfseries,text=alagent!80!black]
      at (9.70,1.55) {Agent\\(LLM)};

    \node[draw=alagent!70,fill=alagent!4,rounded corners=2pt,line width=0.5pt,
          minimum width=\alk*0.85cm,minimum height=\alk*0.75cm] (todo) at (11.75,1.70) {};
    \pic[scale=1.4] at ($(todo.center)+(0.02,0)$) {todo=alagent};
    \node[toolbadge,draw=none,fill=none,minimum width=0pt,anchor=north]
      at ($(todo.south)+(0,-0.03)$) {todo\_list};
    \draw[agentflow,<->] (agent.east |- todo.west) -- (todo.west);

    %% ================================================================ %%
    %% Prompts
    %% ================================================================ %%
    %% --- initial prompt + coverage summary (card grows with text) ---- %%
    \pic at (6.55,3.35) {doc=alprompt};
    \node[alabel,text=alprompt!80!black] (ip) at (6.85,3.35) {initial prompt};
    \draw[alprompt!40,line width=0.4pt]
      ($(ip.east)+(0.25,-0.40)$) -- ($(ip.east)+(0.25,0.40)$);
    \node[anchor=west,inner sep=0pt,text=alprompt!80!black] (cs)
      at ($(ip.east)+(0.50,0.22)$) {coverage summary};
    \fill[alaff!25,rounded corners=1pt]
      ($(cs.south)+(-0.70,-0.12)$) rectangle ++(1.40,-0.18);
    \fill[alcov,rounded corners=1pt]
      ($(cs.south)+(-0.70,-0.12)$) rectangle ++(0.52,-0.18);
    \coordinate (pctop) at (0,3.90);
    \coordinate (pcR) at ($(cs.east)+(0.25,0)$);
    \begin{pgfonlayer}{albg}
      \draw[promptcard] (6.15,2.80) rectangle (pcR |- pctop);
    \end{pgfonlayer}
    \node[prompttag] at (6.15,3.90) {PROMPT};
    \draw[promptflow] (9.35,2.80) -- (agent.north);

    %% --- coverage delta = the next prompt --------------------------- %%
    \draw[promptcard] (9.40,-0.85) rectangle (12.30,0.30);
    \node[prompttag] at (9.40,0.30) {PROMPT};
    \node[text=alprompt!80!black] at (10.85,0.00) {coverage delta};
    \fill[alaff!25,rounded corners=1pt] (9.70,-0.53) rectangle (11.60,-0.35);
    \fill[alcov!45]                    (9.70,-0.53) rectangle (10.40,-0.35);
    \fill[alcov]                       (10.40,-0.53) rectangle (10.90,-0.35);
    \node[font=\alS\bfseries,text=alcov!80!black] at (11.95,-0.44) {+$\Delta$};
    \draw[promptflow] (10.40,-1.20) -- (10.40,-0.85);
    \draw[promptflow] (10.40, 0.30) -- (10.40, 0.80);

    %% ================================================================ %%
    %% Agent -> workspace: tools live on the edges
    %% ================================================================ %%
    \draw[agentflow] (7.95,2.18) -- (4.20,2.18);
    \node[toolbadge] at (6.15,2.18) {write\_file};

    \draw[agentflow] (7.95,0.90) -- (4.20,0.90);
    \matrix (tools) [
        draw=alagent!60, fill=alagent!4, rounded corners=2pt, line width=0.5pt,
        inner sep=0.07cm, column sep=0.06cm, row sep=0.06cm,
        nodes={toolbadge}] at (6.17,0.90) {
      \node{ls};   & \node{hexdump};   \\
      \node{glob}; & \node{strings};   \\
      \node{grep}; & \node{file\_type}; \\
      \node{read\_file}; &            \\
    };
    \node[tooltag] at (tools.north west) {INSPECT};

    %% ================================================================ %%
    %% Agent -> tool: call(path)
    %% ================================================================ %%
    \draw[agentflow] (8.25,0.80) -- (8.25,-1.20);
    \node[draw=alagent!70,fill=white,rounded corners=2pt,line width=0.5pt,
          align=center,inner sep=2.5pt,text=altool] (callarg) at (8.25,-0.42)
      {{\alS\ttfamily\bfseries\textcolor{alagent!80!black}{call(path)}}\\[2pt]
       \tikz[baseline={(0,-0.07)},scale=0.85*\alk]\pic{doc=altool};%
       \hspace{2pt}{\alT\ttfamily
         \begin{tabular}{@{}l@{}}seed\_gen\_\\[-2pt]cycle\_N.py\end{tabular}}};

    %% ================================================================ %%
    %% gain_coverage_report_tool  (container 5.95 .. 12.30 / -5.40 .. -1.20)
    %% ================================================================ %%
    \draw[container] (5.95,-5.40) rectangle (12.30,-1.20);
    \pic at (6.35,-1.52) {gear=altool};
    \node[ctitle,anchor=west] at (6.65,-1.52) {\ttfamily gain\_coverage\_report\_tool};

    \foreach \i/\y in {1/-2.20, 2/-3.08, 3/-3.96, 4/-4.84} {
      \node[steprow] (r\i) at (9.475,\y) {};
      \node[stepnum] (n\i) at (6.40,\y) {2.\i};
    }
    \foreach \a/\b in {n1/n2, n2/n3, n3/n4} \draw[mini] (\a) -- (\b);

    %% [2.1] run + validate
    \pic at (7.15,-2.11) {doc=altool};
    \node[icap] at (7.15,-2.45) {script};
    \draw[mini] (7.36,-2.11) -- (7.60,-2.11);
    \pic at (7.85,-2.11) {gear=altool};
    \node[icap] at (7.85,-2.45) {run};
    \draw[mini] (8.10,-2.11) -- (8.32,-2.11);
    \node[checkbadge,anchor=west] (c1) at (8.36,-2.20) {error};
    \node[checkbadge,right=0.05cm of c1] (c2) {timeout};
    \node[checkbadge,right=0.05cm of c2] (c3) {\#files};
    \node[checkbadge,right=0.05cm of c3] (c4) {size};
    \foreach \c in {c1,c2,c3,c4} \pic at ($(\c.north)+(0,-0.15)$) {check=alcov};

    %% [2.2] persist seeds
    \foreach \x in {7.20,7.46,7.72}
      \pic[scale=0.75] at (\x,-2.99) {doc=alcorpus};
    \node[icap] at (7.46,-3.35) {new seeds};
    \draw[mini] (7.95,-3.03) -- (8.50,-3.03);
    \pic at (8.80,-2.99) {folder=alcorpus};
    \node[icap,text=alcorpus] at (8.80,-3.35) {\ttfamily corpus/};
    \node[icap,anchor=west,font=\alS] at (9.25,-3.03) {persistent};

    %% [2.3] measure coverage on the affected slice
    \pic at (7.25,-3.87) {folder=alcorpus};
    \node[icap] at (7.25,-4.22) {\ttfamily corpus/};
    \draw[mini] (7.52,-3.87) -- (8.02,-3.87);
    \pic at (8.35,-3.87) {chip=altool};
    \node[icap] at (8.35,-4.22) {instr.\ binary};
    \draw[mini] (8.67,-3.87) -- (9.50,-3.87);
    \begin{scope}[shift={(10.10,-3.85)}]
      \fill[alcov,fill opacity=0.22] (-0.17,0) circle (0.22);
      \fill[alaff,fill opacity=0.22] ( 0.17,0) circle (0.22);
      \begin{scope}
        \clip (-0.17,0) circle (0.22);
        \fill[altool,fill opacity=0.45] (0.17,0) circle (0.22);
      \end{scope}
      \draw[alcov,line width=0.5pt] (-0.17,0) circle (0.22);
      \draw[alaff,line width=0.5pt] ( 0.17,0) circle (0.22);
      \node[gmark,text=alcov] at (-0.29,0) {\alT C};
      \node[gmark,text=alaff] at ( 0.29,0) {\alT A};
    \end{scope}
    \node[icap] at (10.10,-4.22) {covered $\cap$ affected};
    \node[icap,anchor=west,font=\alS] at (10.70,-3.85) {slice};

    %% [2.4] update annotation + emit delta
    \node[gmark,text=altool!45] at (7.10,-4.71) {[.]};
    \draw[mini] (7.35,-4.71) -- (7.68,-4.71);
    \node[gmark,text=alcov]     at (7.93,-4.71) {[C]};
    \node[gmark,text=altool!45] at (7.10,-4.97) {[.]};
    \draw[mini] (7.35,-4.97) -- (7.68,-4.97);
    \node[gmark,text=alcov!55]  at (7.93,-4.97) {[P]};
    \draw[mini] (8.22,-4.84) -- (8.55,-4.84);
    \pic at (8.85,-4.77) {folder=altool};
    \node[icap] at (8.85,-5.10) {\ttfamily src/};
    \draw[mini] (9.15,-4.84) -- (10.20,-4.84);
    \node[draw=alprompt,fill=alprompt!10,rounded corners=1.5pt,
          text=alprompt!80!black,font=\alSm\bfseries,
          minimum size=\alk*0.46cm,inner sep=0pt] (delta) at (10.50,-4.84) {$\Delta$};
    \node[icap,anchor=west,text=alprompt!80!black] at (10.80,-4.84) {report};

    %% ================================================================ %%
    %% Tool -> workspace (two lanes; seeds hop over annotate)
    %% ================================================================ %%
    \draw[flow] (5.95,-4.84) -- (5.35,-4.84) -- (5.35,-0.60) -- (4.20,-0.60);
    \node[lbl,rotate=90] at (5.62,-3.65) {annotate};

    \draw[corpusflow,halo] (5.95,-3.08) -- (4.92,-3.08) -- (4.92,-2.00)
      -- (4.20,-2.00);
    \node[corpuslbl,rotate=90] at (4.66,-2.54) {add seeds};

  \end{tikzpicture}
  \caption[Corpus-augmentation loop of the LLM agent]{Corpus-augmentation workspace of the LLM agent}
  \label{fig:agent-loop}
\end{figure}
%
% Required in the preamble (already needed by figures/pipeline.tex):
%   \usepackage{tikz}
%   \usetikzlibrary{arrows.meta,positioning,calc}
%
% Drawing area: see \alk below (tuned for scrbook, fontsize=11pt, paper=a4).
% Do not scale this picture (scale=... shrinks the boxes but not the text).
% All text-bearing boxes size themselves to their content, so the figure
% survives a change of body/typewriter font.

\definecolor{alcorpus}{HTML}{D97706} % corpus spine (orange), match pipeline.tex
\definecolor{alagent}{HTML}{2563EB}  % agent / LLM and its tools
\definecolor{altool}{HTML}{475569}   % tooling and infrastructure
\definecolor{alprompt}{HTML}{7C3AED} % everything the LLM reads as a prompt
\definecolor{alcov}{HTML}{16A34A}    % [C] covered
\definecolor{alaff}{HTML}{DC2626}    % [A] affected
\colorlet{alold}{altool!55}          % older, superseded scripts

% ---- one knob: scales geometry and text together ---------------------------
% 1.125 -> main labels at 9pt = \footnotesize of scrbook (fontsize=11pt),
%          figure is about 14.3cm wide (\textwidth = 14.7cm), 11cm tall.
% 1.0   -> main labels at 8pt, about 12.7cm x 9.8cm.
\def\alk{1.125}
\pgfmathsetmacro\alsN{10*\alk}\pgfmathsetmacro\alsSm{9*\alk}
\pgfmathsetmacro\alsL{12*\alk}\pgfmathsetmacro\alsF{8*\alk}
\pgfmathsetmacro\alsS{7*\alk}\pgfmathsetmacro\alsT{5*\alk}
\def\alN {\fontsize{\alsN }{1.2\dimexpr\alsN  pt\relax}\selectfont}
\def\alSm{\fontsize{\alsSm}{1.2\dimexpr\alsSm pt\relax}\selectfont}
\def\alL {\fontsize{\alsL }{1.2\dimexpr\alsL  pt\relax}\selectfont}
\def\alF {\fontsize{\alsF }{1.2\dimexpr\alsF  pt\relax}\selectfont}
\def\alS {\fontsize{\alsS }{1.2\dimexpr\alsS  pt\relax}\selectfont}
\def\alT {\fontsize{\alsT }{1.2\dimexpr\alsT  pt\relax}\selectfont}

\pgfdeclarelayer{albg}               % card backgrounds drawn behind content
\ifdefined\altoolw\else\newlength{\altoolw}\fi
\settowidth{\altoolw}{\alS\ttfamily file\_type}

\begin{figure}[htbp]
  \centering
  \begin{tikzpicture}[
      scale=\alk,
      >={Latex[length=2.2mm,width=1.6mm]},
      font=\alF,
      %% ---------------- icons (TikZ native, no external fonts) --------- %%
      pics/folder/.style={code={
        \filldraw[draw=#1,fill=#1!14,line width=0.5pt,line join=round]
          (-0.21,-0.16) -- (-0.21,0.15) -- (-0.05,0.15) -- (0.00,0.09)
          -- (0.21,0.09) -- (0.21,-0.16) -- cycle;
        \draw[#1,line width=0.4pt] (-0.21,0.09) -- (0.00,0.09);}},
      pics/doc/.style={code={
        \filldraw[draw=#1,fill=#1!10,line width=0.5pt,line join=round]
          (-0.14,-0.19) -- (-0.14,0.19) -- (0.05,0.19) -- (0.14,0.10)
          -- (0.14,-0.19) -- cycle;
        \draw[#1,line width=0.4pt] (0.05,0.19) -- (0.05,0.10) -- (0.14,0.10);
        \draw[#1!75,line width=0.35pt]
          (-0.08,0.03) -- (0.08,0.03) (-0.08,-0.04) -- (0.08,-0.04)
          (-0.08,-0.11) -- (0.03,-0.11);}},
      pics/chip/.style={code={
        \foreach \y in {-0.08,0.00,0.08} {
          \draw[#1,line width=0.5pt] (-0.16,\y) -- (-0.23,\y);
          \draw[#1,line width=0.5pt] (0.16,\y) -- (0.23,\y);}
        \foreach \x in {-0.08,0.00,0.08} {
          \draw[#1,line width=0.5pt] (\x,0.16) -- (\x,0.23);
          \draw[#1,line width=0.5pt] (\x,-0.16) -- (\x,-0.23);}
        \filldraw[draw=#1,fill=#1!14,line width=0.6pt,rounded corners=1pt]
          (-0.16,-0.16) rectangle (0.16,0.16);
        \draw[#1,line width=0.5pt,rounded corners=0.5pt]
          (-0.07,-0.07) rectangle (0.07,0.07);}},
      pics/gear/.style={code={
        \foreach \a in {0,45,...,315} {
          \filldraw[draw=#1,fill=#1!20,line width=0.4pt]
            (\a-9:0.10) -- (\a-7:0.19) -- (\a+7:0.19) -- (\a+9:0.10) -- cycle;}
        \filldraw[draw=#1,fill=#1!14,line width=0.5pt] (0,0) circle (0.125);
        \filldraw[draw=#1,fill=white,line width=0.4pt] (0,0) circle (0.05);}},
      pics/todo/.style={code={
        \foreach \y/\done in {0.13/1,0.00/1,-0.13/0} {
          \draw[#1,line width=0.4pt,fill=white] (-0.20,\y-0.045) rectangle ++(0.09,0.09);
          \ifnum\done=1
            \draw[#1,line width=0.6pt] (-0.19,\y) -- (-0.16,\y-0.03) -- (-0.10,\y+0.05);
          \fi
          \draw[#1!60,line width=0.9pt,line cap=round] (-0.04,\y) -- (0.18,\y);}}},
      pics/check/.style={code={
        \draw[#1,line width=1.0pt,line cap=round,line join=round]
          (-0.08,0.00) -- (-0.02,-0.06) -- (0.09,0.07);}},
      %% ---------------- boxes ------------------------------------------ %%
      container/.style={
        draw=altool!50, dashed, rounded corners=3pt, line width=0.5pt},
      ctitle/.style={font=\alF\bfseries, text=altool},
      artifact/.style={
        draw=altool!60, fill=altool!6, rounded corners=2pt, line width=0.5pt,
        minimum width=\alk*4.0cm, inner sep=0pt},
      corpusartifact/.style={
        artifact, draw=alcorpus, fill=alcorpus!10, line width=0.7pt},
      alabel/.style={anchor=west, align=left, inner sep=0pt},
      agentbox/.style={
        draw=alagent, fill=alagent!8, rounded corners=3pt, line width=0.8pt,
        minimum width=\alk*2.8cm, minimum height=\alk*1.50cm, inner sep=0pt},
      toolbadge/.style={
        draw=alagent!70, fill=white, rounded corners=1.5pt, line width=0.45pt,
        font=\alS\ttfamily, text=alagent!80!black,
        minimum width=\altoolw+3pt, minimum height=\alk*0.29cm, inner sep=1pt},
      promptcard/.style={
        draw=alprompt, fill=alprompt!6, rounded corners=2pt, line width=0.8pt},
      prompttag/.style={
        fill=alprompt, text=white, font=\alT\bfseries\sffamily,
        rounded corners=1pt, inner sep=1.6pt, anchor=south west},
      finalbox/.style={
        draw=alcorpus, fill=alcorpus!22, rounded corners=2pt, line width=0.8pt,
        minimum width=\alk*4.0cm, minimum height=\alk*0.80cm, inner sep=0pt},
      tooltag/.style={
        fill=alagent, text=white, font=\alT\bfseries\sffamily,
        rounded corners=1pt, inner sep=1.6pt, anchor=south west},
      iotag/.style={
        fill=alcorpus, text=white, font=\alT\bfseries\sffamily,
        rounded corners=1pt, inner sep=1.6pt, anchor=south west},
      steprow/.style={
        draw=altool!45, fill=white, rounded corners=2pt, line width=0.5pt,
        minimum width=\alk*5.25cm, minimum height=\alk*0.80cm, inner sep=0pt},
      stepnum/.style={
        circle, draw=altool, fill=altool, text=white, inner sep=0pt,
        minimum size=\alk*0.46cm, font=\alT\bfseries\sffamily},
      checkbadge/.style={
        draw=alcov!70, fill=alcov!8, rounded corners=1.5pt, line width=0.4pt,
        font=\alT, text=alcov!70!black, minimum width=\alk*0.78cm,
        text height=\alk*0.38cm, text depth=\alk*0.06cm, inner xsep=2pt, inner ysep=1.5pt},
      icap/.style={font=\alT, text=altool, inner sep=0pt},
      gmark/.style={font=\alS\ttfamily\bfseries, inner sep=0pt},
      codeline/.style={line width=2.2pt, line cap=round, draw=altool!30},
      flow/.style={->, draw=altool, line width=0.7pt},
      mini/.style={->, draw=altool!80, line width=0.5pt,
        >={Latex[length=1.3mm,width=1.0mm]}},
      agentflow/.style={->, draw=alagent, line width=0.8pt},
      promptflow/.style={->, draw=alprompt, line width=1.0pt},
      corpusflow/.style={->, draw=alcorpus, line width=1.0pt},
      halo/.style={preaction={draw=white, line width=3.2pt, -}},
      lbl/.style={font=\alS, text=altool, inner sep=1.5pt},
      corpuslbl/.style={lbl, text=alcorpus},
    ]

    \pgfsetlayers{albg,main}

    %% ================================================================ %%
    %% Workspace  (container 0 .. 4.4  /  -3.35 .. 3.30)
    %% ================================================================ %%
    \draw[container] (0,-3.35) rectangle (4.4,3.22);
    \node[ctitle] at (2.2,3.03) {Workspace};

    %% --- generator scripts (one per cycle) --------------------------- %%
    \node[artifact,minimum height=\alk*1.00cm] (script) at (2.2,2.40) {};
    \pic[scale=0.85] at (0.50,2.64) {doc=alold};
    \node[alabel,text=altool!70] at (0.80,2.64) {\ttfamily seed\_gen\_cycle\_1.py};
    \pic[scale=0.85] at (0.50,2.18) {doc=altool};
    \node[alabel] at (0.80,2.18) {\ttfamily seed\_gen\_cycle\_N.py};

    %% --- annotated source (mini listing with gutter markers) -------- %%
    \node[artifact,minimum height=\alk*2.90cm] (src) at (2.2,0.25) {};
    \pic at (0.50,1.40) {folder=altool};
    \node[alabel] at (0.80,1.40)
      {\ttfamily src/};
    \filldraw[draw=altool!30,fill=white,rounded corners=1.5pt,line width=0.4pt]
      (0.35,1.12) rectangle (4.05,-0.42);
    \foreach \y/\m/\c/\x/\len in {
        0.95/[C]/alcov/1.05/2.3,
        0.65/[A]/alaff/1.35/2.3,
        0.35/[P]/alcov!55/1.35/1.5,
        0.05/[A]/alaff/1.65/1.9,
       -0.25/[.]/altool!45/1.05/1.1} {
      \node[gmark,text=\c] at (0.65,\y) {\m};
      \draw[codeline] (\x,\y) -- ++(\len,0);
    }
    \foreach \x/\y/\c/\m/\t in {
        0.40/-0.72/alaff/[A]/affected,  2.25/-0.72/alcov/[C]/covered,
        0.40/-0.98/alcov!55/[P]/partial, 2.25/-0.98/altool!45/[.]/uncovered}
      \node[icap,font=\alS,anchor=west] at (\x,\y)
        {\textcolor{\c}{\ttfamily\bfseries\m}\ \t};

    %% --- corpus: reference + generated seeds ------------------------ %%
    \node[corpusartifact,minimum height=\alk*1.40cm] (corpus) at (2.2,-2.25) {};
    \pic at (0.50,-1.85) {folder=alcorpus};
    \node[alabel] at (0.80,-1.85) {\ttfamily corpus/};
    \foreach \x in {0.65,0.95,1.25} \pic[scale=0.75] at (\x,-2.35) {doc=altool};
    \foreach \x in {2.35,2.65,2.95,3.25} \pic[scale=0.75] at (\x,-2.35) {doc=alcorpus};
    \node[icap] at (0.95,-2.72) {reference};
    \node[icap,text=alcorpus] at (2.80,-2.72) {generated};

    %% --- output ------------------------------------------------------ %%
    \node[finalbox] (final) at (2.2,-4.45) {};
    \node[iotag] at (final.north west) {OUTPUT};
    \pic[scale=1.3] at (0.60,-4.45) {folder=alcorpus};
    \node[alabel,font=\alF\bfseries,text=alcorpus!70!black]
      at (1.05,-4.45) {augmented corpus};
    \draw[corpusflow] (corpus.south) -- (final.north);
    \node[lbl,anchor=west] at (2.35,-3.65) {budget exhausted};

    %% --- input ------------------------------------------------------- %%
    \node[finalbox] (distilled) at (2.2,3.72) {};
    \node[iotag] at (distilled.north west) {INPUT};
    \pic[scale=1.3] at (0.60,3.72) {folder=alcorpus};
    \node[alabel,font=\alF\bfseries,text=alcorpus!70!black]
      at (1.05,3.72) {distilled corpus};
    \draw[corpusflow,line width=1.4pt] (distilled.west) -| (-0.30,-2.25)
      -- (corpus.west);

    %% ================================================================ %%
    %% Agent  (7.75 .. 10.75 / 0.70 .. 2.40) and its todo_list tool
    %% ================================================================ %%
    \node[agentbox] (agent) at (9.35,1.55) {};
    \pic[scale=1.4] at (8.40,1.55) {chip=alagent};
    \node[align=center,font=\alN\bfseries,text=alagent!80!black]
      at (9.70,1.55) {Agent\\(LLM)};

    \node[draw=alagent!70,fill=alagent!4,rounded corners=2pt,line width=0.5pt,
          minimum width=\alk*0.85cm,minimum height=\alk*0.75cm] (todo) at (11.75,1.70) {};
    \pic[scale=1.4] at ($(todo.center)+(0.02,0)$) {todo=alagent};
    \node[toolbadge,draw=none,fill=none,minimum width=0pt,anchor=north]
      at ($(todo.south)+(0,-0.03)$) {todo\_list};
    \draw[agentflow,<->] (agent.east |- todo.west) -- (todo.west);

    %% ================================================================ %%
    %% Prompts
    %% ================================================================ %%
    %% --- initial prompt + coverage summary (card grows with text) ---- %%
    \pic at (6.55,3.35) {doc=alprompt};
    \node[alabel,text=alprompt!80!black] (ip) at (6.85,3.35) {initial prompt};
    \draw[alprompt!40,line width=0.4pt]
      ($(ip.east)+(0.25,-0.40)$) -- ($(ip.east)+(0.25,0.40)$);
    \node[anchor=west,inner sep=0pt,text=alprompt!80!black] (cs)
      at ($(ip.east)+(0.50,0.22)$) {coverage summary};
    \fill[alaff!25,rounded corners=1pt]
      ($(cs.south)+(-0.70,-0.12)$) rectangle ++(1.40,-0.18);
    \fill[alcov,rounded corners=1pt]
      ($(cs.south)+(-0.70,-0.12)$) rectangle ++(0.52,-0.18);
    \coordinate (pctop) at (0,3.90);
    \coordinate (pcR) at ($(cs.east)+(0.25,0)$);
    \begin{pgfonlayer}{albg}
      \draw[promptcard] (6.15,2.80) rectangle (pcR |- pctop);
    \end{pgfonlayer}
    \node[prompttag] at (6.15,3.90) {PROMPT};
    \draw[promptflow] (9.35,2.80) -- (agent.north);

    %% --- coverage delta = the next prompt --------------------------- %%
    \draw[promptcard] (9.40,-0.85) rectangle (12.30,0.30);
    \node[prompttag] at (9.40,0.30) {PROMPT};
    \node[text=alprompt!80!black] at (10.85,0.00) {coverage delta};
    \fill[alaff!25,rounded corners=1pt] (9.70,-0.53) rectangle (11.60,-0.35);
    \fill[alcov!45]                    (9.70,-0.53) rectangle (10.40,-0.35);
    \fill[alcov]                       (10.40,-0.53) rectangle (10.90,-0.35);
    \node[font=\alS\bfseries,text=alcov!80!black] at (11.95,-0.44) {+$\Delta$};
    \draw[promptflow] (10.40,-1.20) -- (10.40,-0.85);
    \draw[promptflow] (10.40, 0.30) -- (10.40, 0.80);

    %% ================================================================ %%
    %% Agent -> workspace: tools live on the edges
    %% ================================================================ %%
    \draw[agentflow] (7.95,2.18) -- (4.20,2.18);
    \node[toolbadge] at (6.15,2.18) {write\_file};

    \draw[agentflow] (7.95,0.90) -- (4.20,0.90);
    \matrix (tools) [
        draw=alagent!60, fill=alagent!4, rounded corners=2pt, line width=0.5pt,
        inner sep=0.07cm, column sep=0.06cm, row sep=0.06cm,
        nodes={toolbadge}] at (6.17,0.90) {
      \node{ls};   & \node{hexdump};   \\
      \node{glob}; & \node{strings};   \\
      \node{grep}; & \node{file\_type}; \\
      \node{read\_file}; &            \\
    };
    \node[tooltag] at (tools.north west) {INSPECT};

    %% ================================================================ %%
    %% Agent -> tool: call(path)
    %% ================================================================ %%
    \draw[agentflow] (8.25,0.80) -- (8.25,-1.20);
    \node[draw=alagent!70,fill=white,rounded corners=2pt,line width=0.5pt,
          align=center,inner sep=2.5pt,text=altool] (callarg) at (8.25,-0.42)
      {{\alS\ttfamily\bfseries\textcolor{alagent!80!black}{call(path)}}\\[2pt]
       \tikz[baseline={(0,-0.07)},scale=0.85*\alk]\pic{doc=altool};%
       \hspace{2pt}{\alT\ttfamily
         \begin{tabular}{@{}l@{}}seed\_gen\_\\[-2pt]cycle\_N.py\end{tabular}}};

    %% ================================================================ %%
    %% gain_coverage_report_tool  (container 5.95 .. 12.30 / -5.40 .. -1.20)
    %% ================================================================ %%
    \draw[container] (5.95,-5.40) rectangle (12.30,-1.20);
    \pic at (6.35,-1.52) {gear=altool};
    \node[ctitle,anchor=west] at (6.65,-1.52) {\ttfamily gain\_coverage\_report\_tool};

    \foreach \i/\y in {1/-2.20, 2/-3.08, 3/-3.96, 4/-4.84} {
      \node[steprow] (r\i) at (9.475,\y) {};
      \node[stepnum] (n\i) at (6.40,\y) {2.\i};
    }
    \foreach \a/\b in {n1/n2, n2/n3, n3/n4} \draw[mini] (\a) -- (\b);

    %% [2.1] run + validate
    \pic at (7.15,-2.11) {doc=altool};
    \node[icap] at (7.15,-2.45) {script};
    \draw[mini] (7.36,-2.11) -- (7.60,-2.11);
    \pic at (7.85,-2.11) {gear=altool};
    \node[icap] at (7.85,-2.45) {run};
    \draw[mini] (8.10,-2.11) -- (8.32,-2.11);
    \node[checkbadge,anchor=west] (c1) at (8.36,-2.20) {error};
    \node[checkbadge,right=0.05cm of c1] (c2) {timeout};
    \node[checkbadge,right=0.05cm of c2] (c3) {\#files};
    \node[checkbadge,right=0.05cm of c3] (c4) {size};
    \foreach \c in {c1,c2,c3,c4} \pic at ($(\c.north)+(0,-0.15)$) {check=alcov};

    %% [2.2] persist seeds
    \foreach \x in {7.20,7.46,7.72}
      \pic[scale=0.75] at (\x,-2.99) {doc=alcorpus};
    \node[icap] at (7.46,-3.35) {new seeds};
    \draw[mini] (7.95,-3.03) -- (8.50,-3.03);
    \pic at (8.80,-2.99) {folder=alcorpus};
    \node[icap,text=alcorpus] at (8.80,-3.35) {\ttfamily corpus/};
    \node[icap,anchor=west,font=\alS] at (9.25,-3.03) {persistent};

    %% [2.3] measure coverage on the affected slice
    \pic at (7.25,-3.87) {folder=alcorpus};
    \node[icap] at (7.25,-4.22) {\ttfamily corpus/};
    \draw[mini] (7.52,-3.87) -- (8.02,-3.87);
    \pic at (8.35,-3.87) {chip=altool};
    \node[icap] at (8.35,-4.22) {instr.\ binary};
    \draw[mini] (8.67,-3.87) -- (9.50,-3.87);
    \begin{scope}[shift={(10.10,-3.85)}]
      \fill[alcov,fill opacity=0.22] (-0.17,0) circle (0.22);
      \fill[alaff,fill opacity=0.22] ( 0.17,0) circle (0.22);
      \begin{scope}
        \clip (-0.17,0) circle (0.22);
        \fill[altool,fill opacity=0.45] (0.17,0) circle (0.22);
      \end{scope}
      \draw[alcov,line width=0.5pt] (-0.17,0) circle (0.22);
      \draw[alaff,line width=0.5pt] ( 0.17,0) circle (0.22);
      \node[gmark,text=alcov] at (-0.29,0) {\alT C};
      \node[gmark,text=alaff] at ( 0.29,0) {\alT A};
    \end{scope}
    \node[icap] at (10.10,-4.22) {covered $\cap$ affected};
    \node[icap,anchor=west,font=\alS] at (10.70,-3.85) {slice};

    %% [2.4] update annotation + emit delta
    \node[gmark,text=altool!45] at (7.10,-4.71) {[.]};
    \draw[mini] (7.35,-4.71) -- (7.68,-4.71);
    \node[gmark,text=alcov]     at (7.93,-4.71) {[C]};
    \node[gmark,text=altool!45] at (7.10,-4.97) {[.]};
    \draw[mini] (7.35,-4.97) -- (7.68,-4.97);
    \node[gmark,text=alcov!55]  at (7.93,-4.97) {[P]};
    \draw[mini] (8.22,-4.84) -- (8.55,-4.84);
    \pic at (8.85,-4.77) {folder=altool};
    \node[icap] at (8.85,-5.10) {\ttfamily src/};
    \draw[mini] (9.15,-4.84) -- (10.20,-4.84);
    \node[draw=alprompt,fill=alprompt!10,rounded corners=1.5pt,
          text=alprompt!80!black,font=\alSm\bfseries,
          minimum size=\alk*0.46cm,inner sep=0pt] (delta) at (10.50,-4.84) {$\Delta$};
    \node[icap,anchor=west,text=alprompt!80!black] at (10.80,-4.84) {report};

    %% ================================================================ %%
    %% Tool -> workspace (two lanes; seeds hop over annotate)
    %% ================================================================ %%
    \draw[flow] (5.95,-4.84) -- (5.35,-4.84) -- (5.35,-0.60) -- (4.20,-0.60);
    \node[lbl,rotate=90] at (5.62,-3.65) {annotate};

    \draw[corpusflow,halo] (5.95,-3.08) -- (4.92,-3.08) -- (4.92,-2.00)
      -- (4.20,-2.00);
    \node[corpuslbl,rotate=90] at (4.66,-2.54) {add seeds};

  \end{tikzpicture}
  \caption[Corpus-augmentation loop of the LLM agent]{Corpus-augmentation workspace of the LLM agent}
  \label{fig:agent-loop}
\end{figure}
%
% Required in the preamble (already needed by figures/pipeline.tex):
%   \usepackage{tikz}
%   \usetikzlibrary{arrows.meta,positioning,calc}
%
% Drawing area: see \alk below (tuned for scrbook, fontsize=11pt, paper=a4).
% Do not scale this picture (scale=... shrinks the boxes but not the text).
% All text-bearing boxes size themselves to their content, so the figure
% survives a change of body/typewriter font.

\definecolor{alcorpus}{HTML}{D97706} % corpus spine (orange), match pipeline.tex
\definecolor{alagent}{HTML}{2563EB}  % agent / LLM and its tools
\definecolor{altool}{HTML}{475569}   % tooling and infrastructure
\definecolor{alprompt}{HTML}{7C3AED} % everything the LLM reads as a prompt
\definecolor{alcov}{HTML}{16A34A}    % [C] covered
\definecolor{alaff}{HTML}{DC2626}    % [A] affected
\colorlet{alold}{altool!55}          % older, superseded scripts

% ---- one knob: scales geometry and text together ---------------------------
% 1.125 -> main labels at 9pt = \footnotesize of scrbook (fontsize=11pt),
%          figure is about 14.3cm wide (\textwidth = 14.7cm), 11cm tall.
% 1.0   -> main labels at 8pt, about 12.7cm x 9.8cm.
\def\alk{1.125}
\pgfmathsetmacro\alsN{10*\alk}\pgfmathsetmacro\alsSm{9*\alk}
\pgfmathsetmacro\alsL{12*\alk}\pgfmathsetmacro\alsF{8*\alk}
\pgfmathsetmacro\alsS{7*\alk}\pgfmathsetmacro\alsT{5*\alk}
\def\alN {\fontsize{\alsN }{1.2\dimexpr\alsN  pt\relax}\selectfont}
\def\alSm{\fontsize{\alsSm}{1.2\dimexpr\alsSm pt\relax}\selectfont}
\def\alL {\fontsize{\alsL }{1.2\dimexpr\alsL  pt\relax}\selectfont}
\def\alF {\fontsize{\alsF }{1.2\dimexpr\alsF  pt\relax}\selectfont}
\def\alS {\fontsize{\alsS }{1.2\dimexpr\alsS  pt\relax}\selectfont}
\def\alT {\fontsize{\alsT }{1.2\dimexpr\alsT  pt\relax}\selectfont}

\pgfdeclarelayer{albg}               % card backgrounds drawn behind content
\ifdefined\altoolw\else\newlength{\altoolw}\fi
\settowidth{\altoolw}{\alS\ttfamily file\_type}

\begin{figure}[htbp]
  \centering
  \begin{tikzpicture}[
      scale=\alk,
      >={Latex[length=2.2mm,width=1.6mm]},
      font=\alF,
      %% ---------------- icons (TikZ native, no external fonts) --------- %%
      pics/folder/.style={code={
        \filldraw[draw=#1,fill=#1!14,line width=0.5pt,line join=round]
          (-0.21,-0.16) -- (-0.21,0.15) -- (-0.05,0.15) -- (0.00,0.09)
          -- (0.21,0.09) -- (0.21,-0.16) -- cycle;
        \draw[#1,line width=0.4pt] (-0.21,0.09) -- (0.00,0.09);}},
      pics/doc/.style={code={
        \filldraw[draw=#1,fill=#1!10,line width=0.5pt,line join=round]
          (-0.14,-0.19) -- (-0.14,0.19) -- (0.05,0.19) -- (0.14,0.10)
          -- (0.14,-0.19) -- cycle;
        \draw[#1,line width=0.4pt] (0.05,0.19) -- (0.05,0.10) -- (0.14,0.10);
        \draw[#1!75,line width=0.35pt]
          (-0.08,0.03) -- (0.08,0.03) (-0.08,-0.04) -- (0.08,-0.04)
          (-0.08,-0.11) -- (0.03,-0.11);}},
      pics/chip/.style={code={
        \foreach \y in {-0.08,0.00,0.08} {
          \draw[#1,line width=0.5pt] (-0.16,\y) -- (-0.23,\y);
          \draw[#1,line width=0.5pt] (0.16,\y) -- (0.23,\y);}
        \foreach \x in {-0.08,0.00,0.08} {
          \draw[#1,line width=0.5pt] (\x,0.16) -- (\x,0.23);
          \draw[#1,line width=0.5pt] (\x,-0.16) -- (\x,-0.23);}
        \filldraw[draw=#1,fill=#1!14,line width=0.6pt,rounded corners=1pt]
          (-0.16,-0.16) rectangle (0.16,0.16);
        \draw[#1,line width=0.5pt,rounded corners=0.5pt]
          (-0.07,-0.07) rectangle (0.07,0.07);}},
      pics/gear/.style={code={
        \foreach \a in {0,45,...,315} {
          \filldraw[draw=#1,fill=#1!20,line width=0.4pt]
            (\a-9:0.10) -- (\a-7:0.19) -- (\a+7:0.19) -- (\a+9:0.10) -- cycle;}
        \filldraw[draw=#1,fill=#1!14,line width=0.5pt] (0,0) circle (0.125);
        \filldraw[draw=#1,fill=white,line width=0.4pt] (0,0) circle (0.05);}},
      pics/todo/.style={code={
        \foreach \y/\done in {0.13/1,0.00/1,-0.13/0} {
          \draw[#1,line width=0.4pt,fill=white] (-0.20,\y-0.045) rectangle ++(0.09,0.09);
          \ifnum\done=1
            \draw[#1,line width=0.6pt] (-0.19,\y) -- (-0.16,\y-0.03) -- (-0.10,\y+0.05);
          \fi
          \draw[#1!60,line width=0.9pt,line cap=round] (-0.04,\y) -- (0.18,\y);}}},
      pics/check/.style={code={
        \draw[#1,line width=1.0pt,line cap=round,line join=round]
          (-0.08,0.00) -- (-0.02,-0.06) -- (0.09,0.07);}},
      %% ---------------- boxes ------------------------------------------ %%
      container/.style={
        draw=altool!50, dashed, rounded corners=3pt, line width=0.5pt},
      ctitle/.style={font=\alF\bfseries, text=altool},
      artifact/.style={
        draw=altool!60, fill=altool!6, rounded corners=2pt, line width=0.5pt,
        minimum width=\alk*4.0cm, inner sep=0pt},
      corpusartifact/.style={
        artifact, draw=alcorpus, fill=alcorpus!10, line width=0.7pt},
      alabel/.style={anchor=west, align=left, inner sep=0pt},
      agentbox/.style={
        draw=alagent, fill=alagent!8, rounded corners=3pt, line width=0.8pt,
        minimum width=\alk*2.8cm, minimum height=\alk*1.50cm, inner sep=0pt},
      toolbadge/.style={
        draw=alagent!70, fill=white, rounded corners=1.5pt, line width=0.45pt,
        font=\alS\ttfamily, text=alagent!80!black,
        minimum width=\altoolw+3pt, minimum height=\alk*0.29cm, inner sep=1pt},
      promptcard/.style={
        draw=alprompt, fill=alprompt!6, rounded corners=2pt, line width=0.8pt},
      prompttag/.style={
        fill=alprompt, text=white, font=\alT\bfseries\sffamily,
        rounded corners=1pt, inner sep=1.6pt, anchor=south west},
      finalbox/.style={
        draw=alcorpus, fill=alcorpus!22, rounded corners=2pt, line width=0.8pt,
        minimum width=\alk*4.0cm, minimum height=\alk*0.80cm, inner sep=0pt},
      tooltag/.style={
        fill=alagent, text=white, font=\alT\bfseries\sffamily,
        rounded corners=1pt, inner sep=1.6pt, anchor=south west},
      iotag/.style={
        fill=alcorpus, text=white, font=\alT\bfseries\sffamily,
        rounded corners=1pt, inner sep=1.6pt, anchor=south west},
      steprow/.style={
        draw=altool!45, fill=white, rounded corners=2pt, line width=0.5pt,
        minimum width=\alk*5.25cm, minimum height=\alk*0.80cm, inner sep=0pt},
      stepnum/.style={
        circle, draw=altool, fill=altool, text=white, inner sep=0pt,
        minimum size=\alk*0.46cm, font=\alT\bfseries\sffamily},
      checkbadge/.style={
        draw=alcov!70, fill=alcov!8, rounded corners=1.5pt, line width=0.4pt,
        font=\alT, text=alcov!70!black, minimum width=\alk*0.78cm,
        text height=\alk*0.38cm, text depth=\alk*0.06cm, inner xsep=2pt, inner ysep=1.5pt},
      icap/.style={font=\alT, text=altool, inner sep=0pt},
      gmark/.style={font=\alS\ttfamily\bfseries, inner sep=0pt},
      codeline/.style={line width=2.2pt, line cap=round, draw=altool!30},
      flow/.style={->, draw=altool, line width=0.7pt},
      mini/.style={->, draw=altool!80, line width=0.5pt,
        >={Latex[length=1.3mm,width=1.0mm]}},
      agentflow/.style={->, draw=alagent, line width=0.8pt},
      promptflow/.style={->, draw=alprompt, line width=1.0pt},
      corpusflow/.style={->, draw=alcorpus, line width=1.0pt},
      halo/.style={preaction={draw=white, line width=3.2pt, -}},
      lbl/.style={font=\alS, text=altool, inner sep=1.5pt},
      corpuslbl/.style={lbl, text=alcorpus},
    ]

    \pgfsetlayers{albg,main}

    %% ================================================================ %%
    %% Workspace  (container 0 .. 4.4  /  -3.35 .. 3.30)
    %% ================================================================ %%
    \draw[container] (0,-3.35) rectangle (4.4,3.22);
    \node[ctitle] at (2.2,3.03) {Workspace};

    %% --- generator scripts (one per cycle) --------------------------- %%
    \node[artifact,minimum height=\alk*1.00cm] (script) at (2.2,2.40) {};
    \pic[scale=0.85] at (0.50,2.64) {doc=alold};
    \node[alabel,text=altool!70] at (0.80,2.64) {\ttfamily seed\_gen\_cycle\_1.py};
    \pic[scale=0.85] at (0.50,2.18) {doc=altool};
    \node[alabel] at (0.80,2.18) {\ttfamily seed\_gen\_cycle\_N.py};

    %% --- annotated source (mini listing with gutter markers) -------- %%
    \node[artifact,minimum height=\alk*2.90cm] (src) at (2.2,0.25) {};
    \pic at (0.50,1.40) {folder=altool};
    \node[alabel] at (0.80,1.40)
      {\ttfamily src/};
    \filldraw[draw=altool!30,fill=white,rounded corners=1.5pt,line width=0.4pt]
      (0.35,1.12) rectangle (4.05,-0.42);
    \foreach \y/\m/\c/\x/\len in {
        0.95/[C]/alcov/1.05/2.3,
        0.65/[A]/alaff/1.35/2.3,
        0.35/[P]/alcov!55/1.35/1.5,
        0.05/[A]/alaff/1.65/1.9,
       -0.25/[.]/altool!45/1.05/1.1} {
      \node[gmark,text=\c] at (0.65,\y) {\m};
      \draw[codeline] (\x,\y) -- ++(\len,0);
    }
    \foreach \x/\y/\c/\m/\t in {
        0.40/-0.72/alaff/[A]/affected,  2.25/-0.72/alcov/[C]/covered,
        0.40/-0.98/alcov!55/[P]/partial, 2.25/-0.98/altool!45/[.]/uncovered}
      \node[icap,font=\alS,anchor=west] at (\x,\y)
        {\textcolor{\c}{\ttfamily\bfseries\m}\ \t};

    %% --- corpus: reference + generated seeds ------------------------ %%
    \node[corpusartifact,minimum height=\alk*1.40cm] (corpus) at (2.2,-2.25) {};
    \pic at (0.50,-1.85) {folder=alcorpus};
    \node[alabel] at (0.80,-1.85) {\ttfamily corpus/};
    \foreach \x in {0.65,0.95,1.25} \pic[scale=0.75] at (\x,-2.35) {doc=altool};
    \foreach \x in {2.35,2.65,2.95,3.25} \pic[scale=0.75] at (\x,-2.35) {doc=alcorpus};
    \node[icap] at (0.95,-2.72) {reference};
    \node[icap,text=alcorpus] at (2.80,-2.72) {generated};

    %% --- output ------------------------------------------------------ %%
    \node[finalbox] (final) at (2.2,-4.45) {};
    \node[iotag] at (final.north west) {OUTPUT};
    \pic[scale=1.3] at (0.60,-4.45) {folder=alcorpus};
    \node[alabel,font=\alF\bfseries,text=alcorpus!70!black]
      at (1.05,-4.45) {augmented corpus};
    \draw[corpusflow] (corpus.south) -- (final.north);
    \node[lbl,anchor=west] at (2.35,-3.65) {budget exhausted};

    %% --- input ------------------------------------------------------- %%
    \node[finalbox] (distilled) at (2.2,3.72) {};
    \node[iotag] at (distilled.north west) {INPUT};
    \pic[scale=1.3] at (0.60,3.72) {folder=alcorpus};
    \node[alabel,font=\alF\bfseries,text=alcorpus!70!black]
      at (1.05,3.72) {distilled corpus};
    \draw[corpusflow,line width=1.4pt] (distilled.west) -| (-0.30,-2.25)
      -- (corpus.west);

    %% ================================================================ %%
    %% Agent  (7.75 .. 10.75 / 0.70 .. 2.40) and its todo_list tool
    %% ================================================================ %%
    \node[agentbox] (agent) at (9.35,1.55) {};
    \pic[scale=1.4] at (8.40,1.55) {chip=alagent};
    \node[align=center,font=\alN\bfseries,text=alagent!80!black]
      at (9.70,1.55) {Agent\\(LLM)};

    \node[draw=alagent!70,fill=alagent!4,rounded corners=2pt,line width=0.5pt,
          minimum width=\alk*0.85cm,minimum height=\alk*0.75cm] (todo) at (11.75,1.70) {};
    \pic[scale=1.4] at ($(todo.center)+(0.02,0)$) {todo=alagent};
    \node[toolbadge,draw=none,fill=none,minimum width=0pt,anchor=north]
      at ($(todo.south)+(0,-0.03)$) {todo\_list};
    \draw[agentflow,<->] (agent.east |- todo.west) -- (todo.west);

    %% ================================================================ %%
    %% Prompts
    %% ================================================================ %%
    %% --- initial prompt + coverage summary (card grows with text) ---- %%
    \pic at (6.55,3.35) {doc=alprompt};
    \node[alabel,text=alprompt!80!black] (ip) at (6.85,3.35) {initial prompt};
    \draw[alprompt!40,line width=0.4pt]
      ($(ip.east)+(0.25,-0.40)$) -- ($(ip.east)+(0.25,0.40)$);
    \node[anchor=west,inner sep=0pt,text=alprompt!80!black] (cs)
      at ($(ip.east)+(0.50,0.22)$) {coverage summary};
    \fill[alaff!25,rounded corners=1pt]
      ($(cs.south)+(-0.70,-0.12)$) rectangle ++(1.40,-0.18);
    \fill[alcov,rounded corners=1pt]
      ($(cs.south)+(-0.70,-0.12)$) rectangle ++(0.52,-0.18);
    \coordinate (pctop) at (0,3.90);
    \coordinate (pcR) at ($(cs.east)+(0.25,0)$);
    \begin{pgfonlayer}{albg}
      \draw[promptcard] (6.15,2.80) rectangle (pcR |- pctop);
    \end{pgfonlayer}
    \node[prompttag] at (6.15,3.90) {PROMPT};
    \draw[promptflow] (9.35,2.80) -- (agent.north);

    %% --- coverage delta = the next prompt --------------------------- %%
    \draw[promptcard] (9.40,-0.85) rectangle (12.30,0.30);
    \node[prompttag] at (9.40,0.30) {PROMPT};
    \node[text=alprompt!80!black] at (10.85,0.00) {coverage delta};
    \fill[alaff!25,rounded corners=1pt] (9.70,-0.53) rectangle (11.60,-0.35);
    \fill[alcov!45]                    (9.70,-0.53) rectangle (10.40,-0.35);
    \fill[alcov]                       (10.40,-0.53) rectangle (10.90,-0.35);
    \node[font=\alS\bfseries,text=alcov!80!black] at (11.95,-0.44) {+$\Delta$};
    \draw[promptflow] (10.40,-1.20) -- (10.40,-0.85);
    \draw[promptflow] (10.40, 0.30) -- (10.40, 0.80);

    %% ================================================================ %%
    %% Agent -> workspace: tools live on the edges
    %% ================================================================ %%
    \draw[agentflow] (7.95,2.18) -- (4.20,2.18);
    \node[toolbadge] at (6.15,2.18) {write\_file};

    \draw[agentflow] (7.95,0.90) -- (4.20,0.90);
    \matrix (tools) [
        draw=alagent!60, fill=alagent!4, rounded corners=2pt, line width=0.5pt,
        inner sep=0.07cm, column sep=0.06cm, row sep=0.06cm,
        nodes={toolbadge}] at (6.17,0.90) {
      \node{ls};   & \node{hexdump};   \\
      \node{glob}; & \node{strings};   \\
      \node{grep}; & \node{file\_type}; \\
      \node{read\_file}; &            \\
    };
    \node[tooltag] at (tools.north west) {INSPECT};

    %% ================================================================ %%
    %% Agent -> tool: call(path)
    %% ================================================================ %%
    \draw[agentflow] (8.25,0.80) -- (8.25,-1.20);
    \node[draw=alagent!70,fill=white,rounded corners=2pt,line width=0.5pt,
          align=center,inner sep=2.5pt,text=altool] (callarg) at (8.25,-0.42)
      {{\alS\ttfamily\bfseries\textcolor{alagent!80!black}{call(path)}}\\[2pt]
       \tikz[baseline={(0,-0.07)},scale=0.85*\alk]\pic{doc=altool};%
       \hspace{2pt}{\alT\ttfamily
         \begin{tabular}{@{}l@{}}seed\_gen\_\\[-2pt]cycle\_N.py\end{tabular}}};

    %% ================================================================ %%
    %% gain_coverage_report_tool  (container 5.95 .. 12.30 / -5.40 .. -1.20)
    %% ================================================================ %%
    \draw[container] (5.95,-5.40) rectangle (12.30,-1.20);
    \pic at (6.35,-1.52) {gear=altool};
    \node[ctitle,anchor=west] at (6.65,-1.52) {\ttfamily gain\_coverage\_report\_tool};

    \foreach \i/\y in {1/-2.20, 2/-3.08, 3/-3.96, 4/-4.84} {
      \node[steprow] (r\i) at (9.475,\y) {};
      \node[stepnum] (n\i) at (6.40,\y) {2.\i};
    }
    \foreach \a/\b in {n1/n2, n2/n3, n3/n4} \draw[mini] (\a) -- (\b);

    %% [2.1] run + validate
    \pic at (7.15,-2.11) {doc=altool};
    \node[icap] at (7.15,-2.45) {script};
    \draw[mini] (7.36,-2.11) -- (7.60,-2.11);
    \pic at (7.85,-2.11) {gear=altool};
    \node[icap] at (7.85,-2.45) {run};
    \draw[mini] (8.10,-2.11) -- (8.32,-2.11);
    \node[checkbadge,anchor=west] (c1) at (8.36,-2.20) {error};
    \node[checkbadge,right=0.05cm of c1] (c2) {timeout};
    \node[checkbadge,right=0.05cm of c2] (c3) {\#files};
    \node[checkbadge,right=0.05cm of c3] (c4) {size};
    \foreach \c in {c1,c2,c3,c4} \pic at ($(\c.north)+(0,-0.15)$) {check=alcov};

    %% [2.2] persist seeds
    \foreach \x in {7.20,7.46,7.72}
      \pic[scale=0.75] at (\x,-2.99) {doc=alcorpus};
    \node[icap] at (7.46,-3.35) {new seeds};
    \draw[mini] (7.95,-3.03) -- (8.50,-3.03);
    \pic at (8.80,-2.99) {folder=alcorpus};
    \node[icap,text=alcorpus] at (8.80,-3.35) {\ttfamily corpus/};
    \node[icap,anchor=west,font=\alS] at (9.25,-3.03) {persistent};

    %% [2.3] measure coverage on the affected slice
    \pic at (7.25,-3.87) {folder=alcorpus};
    \node[icap] at (7.25,-4.22) {\ttfamily corpus/};
    \draw[mini] (7.52,-3.87) -- (8.02,-3.87);
    \pic at (8.35,-3.87) {chip=altool};
    \node[icap] at (8.35,-4.22) {instr.\ binary};
    \draw[mini] (8.67,-3.87) -- (9.50,-3.87);
    \begin{scope}[shift={(10.10,-3.85)}]
      \fill[alcov,fill opacity=0.22] (-0.17,0) circle (0.22);
      \fill[alaff,fill opacity=0.22] ( 0.17,0) circle (0.22);
      \begin{scope}
        \clip (-0.17,0) circle (0.22);
        \fill[altool,fill opacity=0.45] (0.17,0) circle (0.22);
      \end{scope}
      \draw[alcov,line width=0.5pt] (-0.17,0) circle (0.22);
      \draw[alaff,line width=0.5pt] ( 0.17,0) circle (0.22);
      \node[gmark,text=alcov] at (-0.29,0) {\alT C};
      \node[gmark,text=alaff] at ( 0.29,0) {\alT A};
    \end{scope}
    \node[icap] at (10.10,-4.22) {covered $\cap$ affected};
    \node[icap,anchor=west,font=\alS] at (10.70,-3.85) {slice};

    %% [2.4] update annotation + emit delta
    \node[gmark,text=altool!45] at (7.10,-4.71) {[.]};
    \draw[mini] (7.35,-4.71) -- (7.68,-4.71);
    \node[gmark,text=alcov]     at (7.93,-4.71) {[C]};
    \node[gmark,text=altool!45] at (7.10,-4.97) {[.]};
    \draw[mini] (7.35,-4.97) -- (7.68,-4.97);
    \node[gmark,text=alcov!55]  at (7.93,-4.97) {[P]};
    \draw[mini] (8.22,-4.84) -- (8.55,-4.84);
    \pic at (8.85,-4.77) {folder=altool};
    \node[icap] at (8.85,-5.10) {\ttfamily src/};
    \draw[mini] (9.15,-4.84) -- (10.20,-4.84);
    \node[draw=alprompt,fill=alprompt!10,rounded corners=1.5pt,
          text=alprompt!80!black,font=\alSm\bfseries,
          minimum size=\alk*0.46cm,inner sep=0pt] (delta) at (10.50,-4.84) {$\Delta$};
    \node[icap,anchor=west,text=alprompt!80!black] at (10.80,-4.84) {report};

    %% ================================================================ %%
    %% Tool -> workspace (two lanes; seeds hop over annotate)
    %% ================================================================ %%
    \draw[flow] (5.95,-4.84) -- (5.35,-4.84) -- (5.35,-0.60) -- (4.20,-0.60);
    \node[lbl,rotate=90] at (5.62,-3.65) {annotate};

    \draw[corpusflow,halo] (5.95,-3.08) -- (4.92,-3.08) -- (4.92,-2.00)
      -- (4.20,-2.00);
    \node[corpuslbl,rotate=90] at (4.66,-2.54) {add seeds};

  \end{tikzpicture}
  \caption[Corpus-augmentation loop of the LLM agent]{Corpus-augmentation workspace of the LLM agent}
  \label{fig:agent-loop}
\end{figure}

\subsection{Agent Harness}
\label{sec:agent-harness}

The harness is the surrounding program that maintains the conversation state, exposes the available tools to the model, executes the tool calls it emits, and returns their results, as described in \cref{chapter:background}. We use the LangGraph \texttt{deepagents} harness~\cite{DeepAgentsOpen} as our base. It is open-source, well-engineered, straightforward to extend with custom tools, and model-agnostic, allowing the same pipeline to run with different \acp{LLM} without modification. We extend it with the workspace, the toolset, and the coverage-feedback tool described below, while leaving its underlying planning and file-editing behavior unchanged. The base system prompt shipped with \texttt{deepagents} is reproduced in Appendix~\ref{appendix:deepagents_prompt}.

The length of the resulting interaction is not a limiting factor in practice. Modern models provide context windows large enough to accommodate the source excerpts, coverage reports, and generator revisions accumulated over a full augmentation run. Should the context limit nevertheless be approached, the harness summarizes the conversation history and continues from the summary, so that the run is bounded by the budgets described in \cref{sec:agent-termination} rather than by the context window.

\subsection{Workspace}
\label{sec:agent-workspace}

The agent operates on a workspace filesystem rather than on content injected into its prompt. This keeps the amount of material in the context window under the agent's control: it decides which files to open and how much to read from them, rather than receiving the entire program up front. The layout of this workspace is shown in \cref{fig:agent-loop}.

The workspace contains two given artifacts and one that the agent produces. The first is the corpus directory, into which the distilled corpus (see \cref{sec:distilling-corpus}) is placed. It serves a dual role: it is the reference corpus from which the agent infers the expected input format, and it is simultaneously the working corpus to which the generated seeds are added. There is thus no separate output location, and seeds written in one cycle persist and remain visible in all subsequent ones. The second artifact is the source tree of the \ac{SUT}, including the fuzzing harness that defines how an input file is passed into the program under test. The harness is the entry point of every execution, in our case \texttt{LLVMFuzzerTestOneInput}, and therefore determines which parts of the program an input can reach at all; making it available allows the agent to reason about the actual interface it is generating inputs for. The third is the generator script the agent writes, which resides in the workspace root.

The source tree is annotated with the current coverage state, as shown in \cref{fig:annotated-source}. Each line is prefixed with an affectedness marker and a coverage marker, so that the two pieces of information the agent needs are attached directly to the code they describe.
 
% !TeX root = ../main.tex
% figures/annotated_source.tex  --  included via \input{figures/annotated_source.tex}
%
% Required in the preamble:
%   \usepackage{listings}
%   \usepackage{xcolor}
%
% NOTE: the gutter markers are typeset through listings' escape mechanism
% (escapeinside={@}{@}) rather than through `literate`, because `literate`
% inside \lstdefinestyle is unreliable across listings versions and is silently
% overridden by a global \lstset{literate=...} in some preambles. The escape
% character is @, which must therefore not occur in the listed source.

\definecolor{mkA}{HTML}{7C3AED} % [A] affected
\definecolor{mkC}{HTML}{15803D} % [C] covered
\definecolor{mkP}{HTML}{D97706} % [P] partially covered
\definecolor{mkU}{HTML}{B91C1C} % [.] uncovered
\definecolor{srcKw}{HTML}{1D4ED8}
\definecolor{srcCm}{HTML}{6B7280}
\definecolor{srcSt}{HTML}{0F766E}

% coverage-gutter markers; \annB is an empty, exactly marker-wide slot
\newcommand{\annA}{{\ttfamily\color{mkA}[A]}}
\newcommand{\annC}{{\ttfamily\color{mkC}[C]}}
\newcommand{\annP}{{\ttfamily\color{mkP}[P]}}
\newcommand{\annU}{{\ttfamily\color{mkU}[.]}}
\newcommand{\annB}{\hphantom{\ttfamily[A]}}

\lstdefinestyle{annotatedsrc}{
  language=C,
  morekeywords={bool,TRUE,FALSE},
  basicstyle=\fontsize{7}{8}\selectfont\ttfamily,
  keywordstyle=\color{srcKw}\bfseries,
  commentstyle=\color{srcCm}\itshape,
  stringstyle=\color{srcSt},
  columns=fullflexible,
  keepspaces=true,
  showstringspaces=false,
  frame=single,
  rulecolor=\color{black!20},
  framesep=5pt,
  xleftmargin=6pt,
  aboveskip=2pt,
  belowskip=2pt,
  escapeinside={@}{@},
}

\begin{figure}[htbp]
  \centering
  \begin{lstlisting}[style=annotatedsrc]
@\annB\annB@ static void strcpy_url(char *output, const char *url, bool relative)
@\annA\annC@  {
@\annB\annC@    /* we must add this with whitespace-replacing */
@\annB\annC@    bool left = TRUE;
@\annB\annC@    const unsigned char *iptr;
@\annB\annC@    char *optr = output;
@\annB\annC@    const unsigned char *host_sep = (const unsigned char *) url;

@\annB\annC@    if(!relative)
@\annA\annC@      host_sep = (const unsigned char *) find_host_sep(url);

@\annA\annC@    for(iptr = (unsigned char *)url;    /* read from here */
@\annA\annC@        *iptr;         /* until zero byte */
@\annA\annC@        iptr++) {

@\annA\annC@      if(iptr < host_sep) {
@\annA\annC@        *optr++ = *iptr;
@\annB\annC@        continue;
@\annB\annC@      }

@\annA\annC@      switch(*iptr) {
@\annB\annC@      case '?':
@\annA\annC@        left = FALSE;
@\annB\annC@        /* fall through */
@\annB\annC@      default:
@\annA\annP@        if(!ISPRINT(*iptr)) {
@\annA\annU@          snprintf(optr, 4, "%%%02x", *iptr);
@\annB\annU@          optr += 3;
@\annB\annU@        }
@\annB\annB@        else
@\annA\annC@          *optr++=*iptr;
@\annA\annC@        break;
@\annB\annU@      case ' ':
@\annA\annU@        if(left) {
@\annA\annU@          *optr++='%'; /* add a '%' */
@\annB\annU@          *optr++='2'; /* add a '2' */
@\annB\annU@          *optr++='0'; /* add a '0' */
@\annB\annU@        }
@\annB\annB@        else
@\annA\annU@          *optr++='+'; /* add a '+' here */
@\annA\annC@        break;
@\annB\annC@      }
@\annA\annC@    }
@\annA\annC@    *optr = 0; /* zero terminate output buffer */

@\annB\annC@  }
  \end{lstlisting}
  \caption[Coverage-annotated source in the agent workspace]{%
    Excerpt of the coverage-annotated source tree as the agent sees it, here
    \texttt{strcpy\_url} from cURL. The gutter carries an affected marker and
    a coverage marker per line: \annA{} marks membership in the commit-affected
    slice, and \annC, \annP{} and \annU{} denote covered, partially covered, and
    uncovered lines under the current corpus. The agent targets \annA\annU{}
    lines first and \annA\annP{} lines as a fallback. Here the non-printable
    escaping branch and the entire space-encoding branch are unreached, while the
    surrounding loop is covered, which localizes the missing input property to
    URLs containing spaces or non-printable characters.}
  \label{fig:annotated-source}
\end{figure}
 
Here \texttt{[A]} marks a line as belonging to the commit-affected slice computed in \cref{sec:affected-code}, while \texttt{[C]}, \texttt{[P]}, and \texttt{[.]} denote covered, partially covered, and uncovered lines, respectively. Partial coverage arises where a line is reached but not fully evaluated, for instance, under short-circuit evaluation or when a condition holds, but its body is skipped. Coverage is aggregated over the whole corpus rather than attributed to a single execution, so \texttt{[C]} means that some seed covers the line, not that any single seed follows the surrounding path in full. The agent is instructed to target \texttt{[A][.]} lines first and \texttt{[A][P]} lines as a fallback, and to ignore everything else. Placing this information in the source rather than in a separate report means the agent always reads the remaining targets together with the control flow that determines how they are reached. This idea is inspired by SeedMind~\cite{shiHarnessingLargeLanguage2024}, which uses a similar mechanism to combine code with coverage and target signal. The coverage required for the initial annotation is already available from the distillation step, where every seed of the reference corpus was executed on the instrumented binary; the annotation is therefore written once when the workspace is created and updated in place after each cycle.

The agent explores this workspace through the harness's filesystem tools. Directories are traversed with \texttt{ls} and \texttt{glob}, file contents are accessed through \texttt{read\_file}, and \texttt{grep} locates specific identifiers, constants, or format strings across the tree. Because seeds are frequently binary, the agent is additionally given \texttt{hexdump}, \texttt{strings}, and \texttt{file\_type}, which allow it to inspect the byte-level structure of reference seeds, identify headers and length fields, recover embedded textual markers, and determine what kind of input a seed represents. Writing is performed with \texttt{write\_file} and is used only for the generator script.

\subsection{Task Specification}
\label{sec:agent-prompt}

The agent's task is specified through an initial prompt, shown in \cref{fig:agent-prompt}. It states the objective, defines what constitutes a cycle, introduces the coverage-marker legend that gives the annotated source tree its meaning, fixes the persistence contract for generated seeds, describes the two feedback formats, prescribes the per-cycle workflow, and constrains the shape of the generator script. The individual parts are highlighted in the figure to make this structure explicit. A final per-run message binds the cycle budget and the concrete workspace paths and carries the initial coverage summary, which reports the uncovered lines affected by the reference corpus.
 
% !TeX root = ../main.tex
% figures/agent_prompt.tex  --  included via % !TeX root = ../main.tex
% figures/agent_prompt.tex  --  included via % !TeX root = ../main.tex
% figures/agent_prompt.tex  --  included via \input{figures/agent_prompt.tex}
%
% Required in the preamble:
%   \usepackage{listings}
%   \usepackage{xcolor}

\definecolor{ppRole}{HTML}{2563EB}     % role and objective
\definecolor{ppConcept}{HTML}{D97706}  % core concepts: cycle, markers, persistence
\definecolor{ppReport}{HTML}{0F766E}   % feedback formats
\definecolor{ppFlow}{HTML}{7C3AED}     % per-cycle workflow
\definecolor{ppTemplate}{HTML}{475569} % generator script template
\definecolor{ppRun}{HTML}{BE123C}      % per-run instantiation

\lstdefinestyle{promptstyle}{
  basicstyle=\fontsize{6}{7}\selectfont\ttfamily,
  breaklines=true,
  breakautoindent=false,
  breakindent=6pt,
  columns=fullflexible,
  keepspaces=true,
  frame=leftline,
  framerule=1.5pt,
  framesep=4pt,
  xleftmargin=7pt,
  aboveskip=1pt,
  belowskip=5pt,
  showstringspaces=false,
}
% header above each colour-coded block
\newcommand{\promptpart}[2]{%
  \noindent{\color{#1}\scriptsize\bfseries #2}\\[0.5pt]}

\begin{figure}[p]
  \scriptsize
  \centering
  \begin{minipage}[t]{0.485\linewidth}

  \promptpart{ppRole}{Role and objective}
  \begin{lstlisting}[style=promptstyle,rulecolor=\color{ppRole}]
You are an expert in reverse engineering and
coverage-guided fuzzing, specializing in seed corpus
engineering.

Task: Write Python scripts to generate seeds targeting
uncovered affected lines. Your primary objective is to
maximize the number of newly covered affected lines.
Stop when all targetable lines are covered. The
execution environment manages your cycle budget and
will automatically terminate the run when your budget
is spent. Focus entirely on maximizing the coverage
yield of every single tool call.
  \end{lstlisting}

  \promptpart{ppConcept}{Core concepts: cycle, markers, persistence}
  \begin{lstlisting}[style=promptstyle,rulecolor=\color{ppConcept}]
**Cycle**: 1 script written + 1 call to
`gain_coverage_report_tool`. Focus on generating
high-quality, targeted inputs in every single cycle.
Do not stress about tracking or counting cycles; the
environment handles the limits and will halt execution
when exhausted.

**Markers (Coverage Legend)**:
  [A] Affected - depends on the change under test
  [C] Covered (aggregated across runs/seeds; not
      necessarily a single continuous trace)
  [P] Partially covered (e.g., short-circuit
      evaluation, or condition true but body skipped)
  [.] Not covered
Target only [A][.] (first priority) and [A][P]
(fallback). Ignore [A][C] and non-[A] lines.

**Persistence**: Seeds written to `OUTPUT_DIR` persist
across cycles. Your new script only needs to generate
new seeds. Do not edit or modify script files from
previous cycles; always write a new script for the
current cycle.
  \end{lstlisting}

  \promptpart{ppReport}{Feedback formats}
  \begin{lstlisting}[style=promptstyle,rulecolor=\color{ppReport}]
**Initial Coverage Summary**: You will receive a text
summary containing coverage statistics and a mapping
of target lines.
**Tool Delta Output**: The `gain_coverage_report_tool`
returns a text delta showing what changed and what
remains.
  \end{lstlisting}

  \end{minipage}\hfill
  \begin{minipage}[t]{0.485\linewidth}

  \promptpart{ppFlow}{Per-cycle workflow}
  \begin{lstlisting}[style=promptstyle,rulecolor=\color{ppFlow}]
**STRICT WORKFLOW (Log via ToDoList once per cycle,
before coding):**
1. **Target**: Select a set of lines to target from
   the initial `uncovered_affected_lines` list or
   subsequent tool delta `still_uncovered_lines`.
   Prioritize paths that unlock the highest number of
   uncovered affected lines simultaneously. Group
   targets on shared execution paths to maximize
   coverage-per-cycle. Note and skip structurally
   unreachable lines. If no targets remain, stop.
2. **Analyze**: Identify if the SUT expects text
   (JSON, XML, custom) or binary inputs. Trace target
   lines back to the SUT Entry Point
   (`LLVMFuzzerTestOneInput`) to map out parsing
   logic, and study the initial corpus seeds.
3. **Hypothesize**: Design targeted inputs. For text
   targets, use Python's string/formatting tools. For
   binary targets, use Python's `struct` module or
   design structured modifications (swapping headers,
   modifying length fields) of existing seeds to
   reach your targets.
4. **Write**: Write `seed_gen_cycle_N.py` (stdlib
   only) in the root directory (never inside
   `OUTPUT_DIR`) using the `write_file` tool. Use
   targeted generation or structured seed
   modification (no blind random mutation). Ensure
   unique seed filenames.
5. **Execute & Adapt**: Call
   `gain_coverage_report_tool` with the path to the
   new script as the argument. Use the returned
   delta/error to guide your next cycle. If the
   script crashed, your sole priority in the next
   cycle is debugging and fixing the script. Do not
   repeat failed approaches.
  \end{lstlisting}

  \promptpart{ppTemplate}{Generator script template}
  \begin{lstlisting}[style=promptstyle,rulecolor=\color{ppTemplate}]
#!/usr/bin/env python3
import os
OUTPUT_DIR = "corpus"

def generate_targeted_seed():
    """Target: <function/lines> |
       Rationale: <brief description>"""
    filepath = os.path.join(OUTPUT_DIR,
        "seed_cycle_X_description.bin")
    # CONSTRUCT YOUR TARGETED PAYLOAD
    # (Do not leave this empty or generic)
    payload = b""
    with open(filepath, "wb") as f:
        f.write(payload)

if __name__ == "__main__":
    os.makedirs(OUTPUT_DIR, exist_ok=True)
    generators = [generate_targeted_seed]
    for gen in generators:
        gen()
  \end{lstlisting}

  \promptpart{ppRun}{Per-run instantiation}
  \begin{lstlisting}[style=promptstyle,rulecolor=\color{ppRun}]
You have a {cycles}-cycle budget. Treat this as a
strategy heuristic: for small budgets, be aggressive
and target difficult code immediately; for large
budgets, be systematic and build a foundational corpus
first. The environment will handle all termination -
do not waste cognitive effort tracking your current
cycle count.

SUT Entry Point: LLVMFuzzerTestOneInput
Initial Corpus: {initial_corpus_location}
Coverage-Annotated Codebase: {annotated_codebase_location}
Initial Coverage Summary: {initial_coverage_summary}
  \end{lstlisting}

  \end{minipage}

  \caption[Task prompt of the augmentation agent]{Task prompt of the augmentation agent.}
  \label{fig:agent-prompt}
\end{figure}
%
% Required in the preamble:
%   \usepackage{listings}
%   \usepackage{xcolor}

\definecolor{ppRole}{HTML}{2563EB}     % role and objective
\definecolor{ppConcept}{HTML}{D97706}  % core concepts: cycle, markers, persistence
\definecolor{ppReport}{HTML}{0F766E}   % feedback formats
\definecolor{ppFlow}{HTML}{7C3AED}     % per-cycle workflow
\definecolor{ppTemplate}{HTML}{475569} % generator script template
\definecolor{ppRun}{HTML}{BE123C}      % per-run instantiation

\lstdefinestyle{promptstyle}{
  basicstyle=\fontsize{6}{7}\selectfont\ttfamily,
  breaklines=true,
  breakautoindent=false,
  breakindent=6pt,
  columns=fullflexible,
  keepspaces=true,
  frame=leftline,
  framerule=1.5pt,
  framesep=4pt,
  xleftmargin=7pt,
  aboveskip=1pt,
  belowskip=5pt,
  showstringspaces=false,
}
% header above each colour-coded block
\newcommand{\promptpart}[2]{%
  \noindent{\color{#1}\scriptsize\bfseries #2}\\[0.5pt]}

\begin{figure}[p]
  \scriptsize
  \centering
  \begin{minipage}[t]{0.485\linewidth}

  \promptpart{ppRole}{Role and objective}
  \begin{lstlisting}[style=promptstyle,rulecolor=\color{ppRole}]
You are an expert in reverse engineering and
coverage-guided fuzzing, specializing in seed corpus
engineering.

Task: Write Python scripts to generate seeds targeting
uncovered affected lines. Your primary objective is to
maximize the number of newly covered affected lines.
Stop when all targetable lines are covered. The
execution environment manages your cycle budget and
will automatically terminate the run when your budget
is spent. Focus entirely on maximizing the coverage
yield of every single tool call.
  \end{lstlisting}

  \promptpart{ppConcept}{Core concepts: cycle, markers, persistence}
  \begin{lstlisting}[style=promptstyle,rulecolor=\color{ppConcept}]
**Cycle**: 1 script written + 1 call to
`gain_coverage_report_tool`. Focus on generating
high-quality, targeted inputs in every single cycle.
Do not stress about tracking or counting cycles; the
environment handles the limits and will halt execution
when exhausted.

**Markers (Coverage Legend)**:
  [A] Affected - depends on the change under test
  [C] Covered (aggregated across runs/seeds; not
      necessarily a single continuous trace)
  [P] Partially covered (e.g., short-circuit
      evaluation, or condition true but body skipped)
  [.] Not covered
Target only [A][.] (first priority) and [A][P]
(fallback). Ignore [A][C] and non-[A] lines.

**Persistence**: Seeds written to `OUTPUT_DIR` persist
across cycles. Your new script only needs to generate
new seeds. Do not edit or modify script files from
previous cycles; always write a new script for the
current cycle.
  \end{lstlisting}

  \promptpart{ppReport}{Feedback formats}
  \begin{lstlisting}[style=promptstyle,rulecolor=\color{ppReport}]
**Initial Coverage Summary**: You will receive a text
summary containing coverage statistics and a mapping
of target lines.
**Tool Delta Output**: The `gain_coverage_report_tool`
returns a text delta showing what changed and what
remains.
  \end{lstlisting}

  \end{minipage}\hfill
  \begin{minipage}[t]{0.485\linewidth}

  \promptpart{ppFlow}{Per-cycle workflow}
  \begin{lstlisting}[style=promptstyle,rulecolor=\color{ppFlow}]
**STRICT WORKFLOW (Log via ToDoList once per cycle,
before coding):**
1. **Target**: Select a set of lines to target from
   the initial `uncovered_affected_lines` list or
   subsequent tool delta `still_uncovered_lines`.
   Prioritize paths that unlock the highest number of
   uncovered affected lines simultaneously. Group
   targets on shared execution paths to maximize
   coverage-per-cycle. Note and skip structurally
   unreachable lines. If no targets remain, stop.
2. **Analyze**: Identify if the SUT expects text
   (JSON, XML, custom) or binary inputs. Trace target
   lines back to the SUT Entry Point
   (`LLVMFuzzerTestOneInput`) to map out parsing
   logic, and study the initial corpus seeds.
3. **Hypothesize**: Design targeted inputs. For text
   targets, use Python's string/formatting tools. For
   binary targets, use Python's `struct` module or
   design structured modifications (swapping headers,
   modifying length fields) of existing seeds to
   reach your targets.
4. **Write**: Write `seed_gen_cycle_N.py` (stdlib
   only) in the root directory (never inside
   `OUTPUT_DIR`) using the `write_file` tool. Use
   targeted generation or structured seed
   modification (no blind random mutation). Ensure
   unique seed filenames.
5. **Execute & Adapt**: Call
   `gain_coverage_report_tool` with the path to the
   new script as the argument. Use the returned
   delta/error to guide your next cycle. If the
   script crashed, your sole priority in the next
   cycle is debugging and fixing the script. Do not
   repeat failed approaches.
  \end{lstlisting}

  \promptpart{ppTemplate}{Generator script template}
  \begin{lstlisting}[style=promptstyle,rulecolor=\color{ppTemplate}]
#!/usr/bin/env python3
import os
OUTPUT_DIR = "corpus"

def generate_targeted_seed():
    """Target: <function/lines> |
       Rationale: <brief description>"""
    filepath = os.path.join(OUTPUT_DIR,
        "seed_cycle_X_description.bin")
    # CONSTRUCT YOUR TARGETED PAYLOAD
    # (Do not leave this empty or generic)
    payload = b""
    with open(filepath, "wb") as f:
        f.write(payload)

if __name__ == "__main__":
    os.makedirs(OUTPUT_DIR, exist_ok=True)
    generators = [generate_targeted_seed]
    for gen in generators:
        gen()
  \end{lstlisting}

  \promptpart{ppRun}{Per-run instantiation}
  \begin{lstlisting}[style=promptstyle,rulecolor=\color{ppRun}]
You have a {cycles}-cycle budget. Treat this as a
strategy heuristic: for small budgets, be aggressive
and target difficult code immediately; for large
budgets, be systematic and build a foundational corpus
first. The environment will handle all termination -
do not waste cognitive effort tracking your current
cycle count.

SUT Entry Point: LLVMFuzzerTestOneInput
Initial Corpus: {initial_corpus_location}
Coverage-Annotated Codebase: {annotated_codebase_location}
Initial Coverage Summary: {initial_coverage_summary}
  \end{lstlisting}

  \end{minipage}

  \caption[Task prompt of the augmentation agent]{Task prompt of the augmentation agent.}
  \label{fig:agent-prompt}
\end{figure}
%
% Required in the preamble:
%   \usepackage{listings}
%   \usepackage{xcolor}

\definecolor{ppRole}{HTML}{2563EB}     % role and objective
\definecolor{ppConcept}{HTML}{D97706}  % core concepts: cycle, markers, persistence
\definecolor{ppReport}{HTML}{0F766E}   % feedback formats
\definecolor{ppFlow}{HTML}{7C3AED}     % per-cycle workflow
\definecolor{ppTemplate}{HTML}{475569} % generator script template
\definecolor{ppRun}{HTML}{BE123C}      % per-run instantiation

\lstdefinestyle{promptstyle}{
  basicstyle=\fontsize{6}{7}\selectfont\ttfamily,
  breaklines=true,
  breakautoindent=false,
  breakindent=6pt,
  columns=fullflexible,
  keepspaces=true,
  frame=leftline,
  framerule=1.5pt,
  framesep=4pt,
  xleftmargin=7pt,
  aboveskip=1pt,
  belowskip=5pt,
  showstringspaces=false,
}
% header above each colour-coded block
\newcommand{\promptpart}[2]{%
  \noindent{\color{#1}\scriptsize\bfseries #2}\\[0.5pt]}

\begin{figure}[p]
  \scriptsize
  \centering
  \begin{minipage}[t]{0.485\linewidth}

  \promptpart{ppRole}{Role and objective}
  \begin{lstlisting}[style=promptstyle,rulecolor=\color{ppRole}]
You are an expert in reverse engineering and
coverage-guided fuzzing, specializing in seed corpus
engineering.

Task: Write Python scripts to generate seeds targeting
uncovered affected lines. Your primary objective is to
maximize the number of newly covered affected lines.
Stop when all targetable lines are covered. The
execution environment manages your cycle budget and
will automatically terminate the run when your budget
is spent. Focus entirely on maximizing the coverage
yield of every single tool call.
  \end{lstlisting}

  \promptpart{ppConcept}{Core concepts: cycle, markers, persistence}
  \begin{lstlisting}[style=promptstyle,rulecolor=\color{ppConcept}]
**Cycle**: 1 script written + 1 call to
`gain_coverage_report_tool`. Focus on generating
high-quality, targeted inputs in every single cycle.
Do not stress about tracking or counting cycles; the
environment handles the limits and will halt execution
when exhausted.

**Markers (Coverage Legend)**:
  [A] Affected - depends on the change under test
  [C] Covered (aggregated across runs/seeds; not
      necessarily a single continuous trace)
  [P] Partially covered (e.g., short-circuit
      evaluation, or condition true but body skipped)
  [.] Not covered
Target only [A][.] (first priority) and [A][P]
(fallback). Ignore [A][C] and non-[A] lines.

**Persistence**: Seeds written to `OUTPUT_DIR` persist
across cycles. Your new script only needs to generate
new seeds. Do not edit or modify script files from
previous cycles; always write a new script for the
current cycle.
  \end{lstlisting}

  \promptpart{ppReport}{Feedback formats}
  \begin{lstlisting}[style=promptstyle,rulecolor=\color{ppReport}]
**Initial Coverage Summary**: You will receive a text
summary containing coverage statistics and a mapping
of target lines.
**Tool Delta Output**: The `gain_coverage_report_tool`
returns a text delta showing what changed and what
remains.
  \end{lstlisting}

  \end{minipage}\hfill
  \begin{minipage}[t]{0.485\linewidth}

  \promptpart{ppFlow}{Per-cycle workflow}
  \begin{lstlisting}[style=promptstyle,rulecolor=\color{ppFlow}]
**STRICT WORKFLOW (Log via ToDoList once per cycle,
before coding):**
1. **Target**: Select a set of lines to target from
   the initial `uncovered_affected_lines` list or
   subsequent tool delta `still_uncovered_lines`.
   Prioritize paths that unlock the highest number of
   uncovered affected lines simultaneously. Group
   targets on shared execution paths to maximize
   coverage-per-cycle. Note and skip structurally
   unreachable lines. If no targets remain, stop.
2. **Analyze**: Identify if the SUT expects text
   (JSON, XML, custom) or binary inputs. Trace target
   lines back to the SUT Entry Point
   (`LLVMFuzzerTestOneInput`) to map out parsing
   logic, and study the initial corpus seeds.
3. **Hypothesize**: Design targeted inputs. For text
   targets, use Python's string/formatting tools. For
   binary targets, use Python's `struct` module or
   design structured modifications (swapping headers,
   modifying length fields) of existing seeds to
   reach your targets.
4. **Write**: Write `seed_gen_cycle_N.py` (stdlib
   only) in the root directory (never inside
   `OUTPUT_DIR`) using the `write_file` tool. Use
   targeted generation or structured seed
   modification (no blind random mutation). Ensure
   unique seed filenames.
5. **Execute & Adapt**: Call
   `gain_coverage_report_tool` with the path to the
   new script as the argument. Use the returned
   delta/error to guide your next cycle. If the
   script crashed, your sole priority in the next
   cycle is debugging and fixing the script. Do not
   repeat failed approaches.
  \end{lstlisting}

  \promptpart{ppTemplate}{Generator script template}
  \begin{lstlisting}[style=promptstyle,rulecolor=\color{ppTemplate}]
#!/usr/bin/env python3
import os
OUTPUT_DIR = "corpus"

def generate_targeted_seed():
    """Target: <function/lines> |
       Rationale: <brief description>"""
    filepath = os.path.join(OUTPUT_DIR,
        "seed_cycle_X_description.bin")
    # CONSTRUCT YOUR TARGETED PAYLOAD
    # (Do not leave this empty or generic)
    payload = b""
    with open(filepath, "wb") as f:
        f.write(payload)

if __name__ == "__main__":
    os.makedirs(OUTPUT_DIR, exist_ok=True)
    generators = [generate_targeted_seed]
    for gen in generators:
        gen()
  \end{lstlisting}

  \promptpart{ppRun}{Per-run instantiation}
  \begin{lstlisting}[style=promptstyle,rulecolor=\color{ppRun}]
You have a {cycles}-cycle budget. Treat this as a
strategy heuristic: for small budgets, be aggressive
and target difficult code immediately; for large
budgets, be systematic and build a foundational corpus
first. The environment will handle all termination -
do not waste cognitive effort tracking your current
cycle count.

SUT Entry Point: LLVMFuzzerTestOneInput
Initial Corpus: {initial_corpus_location}
Coverage-Annotated Codebase: {annotated_codebase_location}
Initial Coverage Summary: {initial_coverage_summary}
  \end{lstlisting}

  \end{minipage}

  \caption[Task prompt of the augmentation agent]{Task prompt of the augmentation agent.}
  \label{fig:agent-prompt}
\end{figure}

The prompt was not derived analytically. Its wording and structure are the result of an iterative pilot phase: we ran the agent on the open-source projects \emph{readstat} and \emph{libhtp}, inspected the resulting transcripts, and revised the specification whenever a systematic deviation from the intended workflow appeared; most notably, trying to blindly generate inputs without inspecting the existing corpus and codebase, generating non-harness-compatible Python seed generation scripts, or focusing on only covering a small hard to reach fraction of code. Iteration stopped once no such deviation recurred across many consecutive pilot runs. The prompt was then frozen and used verbatim for every run reported in \cref{chapter:evaluation}; no per-project or per-arm adaptation was performed. We make no claim that this formulation is optimal, only that it is fixed and identical across arms, so that any observed difference is attributable to the starting corpus rather than to the instructions.
 
Beyond this specification, the prompt prescribes no strategy: it does not tell the agent which files to read, in which order to investigate the program, or how to construct inputs. Determining this is the task itself. What the workflow does constrain is the discipline around that investigation. The agent is required to select and record its targets before writing code, to prefer targets that lie on shared execution paths so that a single cycle can unlock several affected lines at once, to note and skip lines it judges structurally unreachable rather than spending further cycles on them, and to prioritize debugging over new construction whenever its previous script failed to execute.

\subsection{Seed Generation}
\label{sec:agent-generation}

Given this environment, the agent follows an iterative corpus-augmentation procedure. It first examines the source tree and the corpus to identify the expected input format, the parsing and validation logic that inputs must survive, and the conditions that guard the paths leading to the uncovered affected lines. It then writes a Python script that produces new seed files in the corpus directory.
 
The agent, therefore, does not emit seeds directly, but a program that emits them. This indirection follows the design adopted by prior \ac{LLM}-based seed-generation work~\cite{zhuSeedSmithLLMDrivenSeed2026,xuISC4DGFEnhancingDirected2024,fengFuzzingRandomnessReasoning2025} and is motivated by two considerations. First, it raises the level of abstraction: emitting raw bytes or hexadecimal literals token by token is error-prone and scales poorly to inputs with dependent fields, whereas a script can express structure explicitly, compute length and offset fields from the data they describe, and emit families of related inputs from a single parameterized construction. Second, it grants access to Python's standard library, so that checksums, compression, encoding, and arithmetic transformations, which are frequently required to pass format validation, can be computed correctly rather than approximated by the model. Generators are restricted to the standard library, both to keep the execution environment reproducible and to prevent the agent from delegating input construction to an existing format library rather than deriving the format from the program.
 
Each cycle produces a new script, named \texttt{seed\_gen\_cycle\_N.py}, rather than an edit of the previous one. This follows from the persistence of the corpus directory: because the seeds of earlier cycles are already on disk, a script never needs to be re-run, and revising it would risk regressing constructions that previously worked. Each script is therefore responsible only for the seeds of its own cycle, and the sequence of scripts additionally forms a readable record of how the agent's understanding of the format developed over the run.
 
When the agent considers its generator ready, it ends the cycle by invoking the coverage report tool with the path to the script.

\subsection{Coverage Feedback}
\label{sec:agent-feedback}

The \texttt{gain\_coverage\_report\_tool} is the only means by which the agent affects the corpus and the only source of feedback it receives. It performs the following steps, corresponding to sub-steps 2.1--2.4 in \cref{fig:agent-loop}. 
 
First, the generator script is executed in an isolated working directory under a wall-clock timeout. The execution fails if the script contains a syntax error, raises an exception, exceeds the timeout, produces no files, produces more files than the configured limit, or produces files exceeding the configured maximum size. In all of these cases, the tool returns a failure message describing the cause, including the interpreter traceback when available, and the corpus remains unchanged. The file count and size limits prevent the agent from replacing targeted construction with bulk generation, which would both defeat the corpus-size objective of \cref{sec:design-goals} and shift the burden back onto the fuzzer.
 
Second, on success, the generated files are added to the corpus. We retain all generated seeds, including those that contribute no additional affected-code coverage. A seed that fails to reach affected code may still be structurally valid and lie closer to it in the mutation space than anything already present, and thus remain a useful starting point for the fuzzer even though it is not credited by the coverage signal.
 
Third, the resulting corpus is executed on the coverage-instrumented binary, and its line coverage is aggregated across all seeds and intersected with the affected slice. The annotated source tree in the workspace is rewritten in place to reflect this new state, so that the agent's view of the program remains consistent with the corpus it has produced.
 
Finally, a coverage report is returned to the agent, as shown in Figure~\cref{fig:coverage-report}. Rather than restating the full coverage state, it is formulated as a delta: it reports which affected lines the cycle newly covered and which remain uncovered, together with the resulting \ac{CACC}. The list of still-uncovered lines is the input to the target-selection step of the next cycle.

% !TeX root = ../main.tex
% figures/coverage_report.tex  --  included via \input{figures/coverage_report.tex}
%
% Required in the preamble:
%   \usepackage{listings}
%   \usepackage{xcolor}

\definecolor{rptGain}{HTML}{15803D} % newly covered
\definecolor{rptLeft}{HTML}{B91C1C} % still uncovered
\definecolor{rptLoss}{HTML}{475569} % newly uncovered (regression guard)
\definecolor{rptFade}{HTML}{9CA3AF} % elisions

\newcommand{\rptGainHead}[1]{{\color{rptGain}\bfseries #1}}
\newcommand{\rptLeftHead}[1]{{\color{rptLeft}\bfseries #1}}
\newcommand{\rptLossHead}[1]{{\color{rptLoss}\bfseries #1}}
\newcommand{\rptElide}[1]{{\color{rptFade}\itshape #1}}

\lstdefinestyle{coveragereport}{
  basicstyle=\fontsize{7}{8.4}\selectfont\ttfamily,
  columns=fullflexible,
  keepspaces=true,
  showstringspaces=false,
  frame=single,
  rulecolor=\color{black!20},
  framesep=5pt,
  xleftmargin=6pt,
  aboveskip=2pt,
  belowskip=2pt,
  escapeinside={(*}{*)},
}

\begin{figure}[htbp]
  \centering
  \begin{lstlisting}[style=coveragereport]
(*\rptGainHead{Newly covered affected lines: 5595}*)
  /src/curl/lib/../lib/sigpipe.h     43, 48, 57, 65, 68-69
  /src/curl/lib/asyn-thread.c        93, 112, 116-117, 126, (*\rptElide{[...]}*)
  /src/curl/lib/base64.c             175, 198, 207, 213, 216-219, (*\rptElide{[...]}*)
  (*\rptElide{[...]}*)

(*\rptLeftHead{Still uncovered affected lines: 11039}*)
  /src/curl/lib/url.c                808/1336 affected lines uncovered ( 60%)
  /src/curl/lib/ftp.c                766/1170 affected lines uncovered ( 65%)
  /src/curl/lib/vtls/openssl.c       679/ 806 affected lines uncovered ( 84%)
  /src/curl/lib/setopt.c             594/ 683 affected lines uncovered ( 87%)
  /src/curl/lib/http.c               556/1196 affected lines uncovered ( 46%)
  /src/curl/lib/multi.c              481/ 868 affected lines uncovered ( 55%)
  /src/curl/lib/http2.c              450/ 470 affected lines uncovered ( 96%)
  /src/curl/lib/telnet.c             396/ 396 affected lines uncovered (100%)
  /src/curl/lib/ftplistparser.c      345/ 349 affected lines uncovered ( 99%)
  /src/curl/lib/transfer.c           330/ 581 affected lines uncovered ( 57%)
  (*\rptElide{[...]}*)
  (*\rptElide{... 42 more files, 488 more uncovered affected lines}*)

(*\rptLossHead{Newly uncovered affected lines: 0}*)
  \end{lstlisting}
  \caption[Coverage delta report returned to the agent]{%
    Coverage delta returned by \texttt{gain\_coverage\_report\_tool} after a
    cycle, here for cURL. The report is a delta rather than a full coverage
    dump: it names the affected lines the cycle newly reached, then ranks the
    files that still contain uncovered affected lines by absolute count, giving
    the uncovered fraction per file so that the agent can distinguish files it
    has barely entered from files it has almost exhausted. The trailing count
    closes the list without enumerating the long tail. Elisions marked
    \rptElide{[...]} abbreviate the report for presentation.}
  \label{fig:coverage-report}
\end{figure}
 
Both the report and the source annotation describe the corpus as a whole rather than individual seeds. This is a deliberate trade-off in favor of context efficiency: per-seed reports would grow with every cycle and quickly dominate the context window, whereas an aggregate report has a constant size regardless of how many seeds the corpus contains. The cost is a coarser signal: the agent learns whether its most recent batch increased coverage, but not which specific input caused the increase.

\subsection{Termination}
\label{sec:agent-termination}

The agent repeats this cycle, revising its generator in response to each report, until one of two budgets is exhausted. The first is a cycle budget that bounds the number of \texttt{get\_coverage\_report} invocations, analogous to the epoch count used during model training. The second is a token budget with soft and hard limits: upon reaching the soft limit, the agent may finalize its current generator and terminate the run; upon reaching the hard limit, the run terminates unconditionally. The soft limit exists so that a run ending on budget still ends after a completed cycle rather than mid-revision. If the agent ends its iteration before the budget is reached, it is prompted to continue executing until the budget is exhausted.

The working corpus as of the last successful cycle constitutes the final augmented corpus $C^*$, which serves as the starting corpus for the regression-fuzzing campaign.